\documentclass[pdflatex,sn-mathphys-num]{sn-jnl}

\usepackage{graphicx}%
\usepackage{multirow}%
\usepackage{amsmath,amssymb,amsfonts}%
\usepackage{amsthm}%
\usepackage{mathrsfs}%
\usepackage[title]{appendix}%
\usepackage{xcolor}%
\usepackage{soul}
\usepackage{subcaption}
\usepackage{booktabs}


\begin{document}

\title[Structure-guided prioritization of divergent candidates in E. coli]{Structure-guided prioritization of divergent virulence and resistance candidates in an open pangenome clinical \textit{Escherichia coli} isolate}

\author*[1]{\fnm{Sarra} \sur{Benmoumou-Hosni}}\email{sarra\_ben@outlook.fr}

\author[1]{\fnm{Atika} \sur{Meklat}}

\author[2]{\fnm{Ikram} \sur{Haleche}}

\affil*[1]{\orgdiv{Laboratoire de Biologie des Systèmes Microbiens}, \orgname{École Normale Supérieure de Kouba}, \orgaddress{\city{Algiers}, \postcode{16000}, \country{Algeria}}}
\affil[2]{\orgname{Université Saad Dahlab Blida 1}, \orgaddress{\street{B.P. 270, Route de Soumâa}, \city{Blida}, \postcode{09000}, \country{Algeria}}}

\abstract{Determining the biological function of hypothetical proteins remains a major bottleneck in bacterial genomics. The accessory genome of clinical pathogens conceals sequence-divergent antimicrobial resistance (AMR) and virulence determinants that escape annotation by standard sequence-alignment tools. We present an integrated, structure-guided computational prioritization pipeline combining pangenomic singleton filtering, high-throughput structural prediction (ESMFold and AlphaFold3), structural homology searches (Foldseek/TM-align), protein language model (PLM) embeddings (ESM-2 650M), and explicit-solvent molecular dynamics (MD) simulations (OpenMM). We applied this pipeline to the clinical multidrug-resistant (MDR) extraintestinal pathogenic \textit{Escherichia coli} (ExPEC) sequence type 354 (ST354) isolate QA5221. Pangenome orthology clustering across 32 lineage genomes isolated 142 private singletons; sequence length and domain filters identified 22 prioritized candidates. The structure-guided pipeline prioritized four candidates representing highly divergent homologs within known virulence and resistance superfamilies: (i) GNAT\_KA27 (\textit{KNGPFPPJ\_02769}), a putative GCN5-related N-acetyltransferase (GNAT) showing a mean predicted aminoglycoside binding free energy of $\Delta G_{\text{bind}} = -22.91$~kcal/mol (averaged across kanamycin: $-23.90$, gentamicin: $-22.84$, and amikacin: $-21.99$~kcal/mol) via computational MM-GBSA and 127-ns MD fold stability (RMSD = $2.73 \pm 0.20$ \AA); (ii) Ehly\_61 (\textit{KNGPFPPJ\_00061}), a putative enterohemolysin showing high structural stability in explicit phospholipid membrane bilayer MD simulation, maintaining a stable fold ($5.62 \pm 1.04$ \AA\ RMSD, $R_g = 27.25 \pm 0.21$ \AA) over 90~ns; (iii) OAgP\_161 (\textit{KNGPFPPJ\_03161}), a putative O-antigen polymerase (Wzy) located within the lipopolysaccharide gene cluster showing stable transmembrane fold topology ($3.10 \pm 0.24$ \AA\ RMSD); and (iv) OAT\_371 (\textit{KNGPFPPJ\_04371}), a plasmid-associated outer-membrane autotransporter candidate representing the weakest evidence tier among the prioritized targets. All four genes are ultra-rare singletons, feature low GC content (29--48\%), and cluster distinctly within the ESM-2 latent space. Family-level phylogenetic analyses confirm these candidates form independent, deeply branching clades distinct from characterized homologs. This computational framework enables the systematic prioritization of sequence-divergent accessory genes, flagging high-priority candidates for subsequent experimental characterization.}

\keywords{Escherichia coli, antimicrobial resistance, pangenome-to-structure, molecular dynamics, protein language models, extraintestinal pathogenic \textit{Escherichia coli} (ExPEC)}

\maketitle


\section{Introduction}\label{sec:introduction}
Antimicrobial resistance (AMR) represents one of the most critical threats to global public health in the 21st century, causing an estimated 1.27 million deaths annually \cite{murray2022global}. Extraintestinal pathogenic \textit{Escherichia coli} (ExPEC) lineages, particularly multidrug-resistant (MDR) clonal groups like Sequence Type 354 (ST354), are primary drivers of this crisis, causing a wide spectrum of severe diseases such as urinary tract infections, bacteremia, and neonatal meningitis \cite{marin2020expec,fuga2024hybrid}. In these pathogenic lineages, treatment failure is frequently driven by a combination of acquired antibiotic resistance genes (ARGs) and specific chromosomal mutations in drug targets, such as those in the quinolone resistance-determining region (QRDR), which together restrict therapeutic options and exacerbate clinical outcomes \cite{bortolaia2020resfinder,feldgarden2019val}.

The genomic landscape of \textit{E. coli} is characterized by an open pangenome, which continues to expand as new genomes are sequenced, reflecting a highly dynamic evolutionary trajectory \cite{touchon2009organised,ghalayini2022distinct}. While core genes are conserved across all isolates to maintain basic cellular functions, the accessory genome consists of shell and cloud partitions that represent a vast, strain-specific reservoir of genetic diversity, containing private singletons or ``ORFans'' that are unique to individual strains. This accessory fraction acts as a primary vehicle for genomic adaptation, frequently accumulating novel virulence and resistance determinants acquired through horizontal gene transfer (HGT) from diverse environmental and clinical donor pools \cite{partridge2018mobile,silva2022evolutionary}.

Mobile genetic elements (MGEs) such as insertion sequences (ISs), integrons, and plasmids serve as the primary drivers of this genomic plasticity \cite{partridge2018mobile,gillings2015class}. These elements act as natural gene-capture systems, facilitating the transposition, recombination, and dissemination of multidrug resistance cassettes and virulence factors across diverse bacterial populations \cite{carattoli2013plasmid,robertson2020mob}. Furthermore, insertion sequences can disrupt ancestral regulatory regions or mobilize adjacent chromosomal segments, giving rise to novel gene arrangements and atypical resistance phenotypes that enable rapid adaptation to host-selective and antibiotic pressures \cite{xie2017isescan,bouvier2022integronfinder}.

Despite the abundance of genomic data, a substantial fraction of predicted coding sequences in draft bacterial assemblies—often exceeding 30\%—are classified as hypothetical proteins due to the limitations of standard annotation pipelines. Conventional sequence-alignment tools, such as BLASTp \cite{altschul1990basic} and profile Hidden Markov Models (HMMs) \cite{eddy2011accelerated,mistry2021pfam}, rely on primary sequence identity and fail to annotate highly divergent homologs that share less than 20--25\% amino acid identity with characterized templates. These proteins occupy the sequence ``twilight zone'' of homology, where sequence similarity is lost but biological function remains conserved, highlighting the critical need for sequence-independent annotation tools that leverage structural topology and machine learning \cite{mirdita2022foldseek,zhang2005tm}.

The recent development of high-accuracy protein structure prediction algorithms, such as AlphaFold and ESMFold, has transformed structural biology by resolving tertiary configurations directly from primary sequences with atomic-level precision \cite{jumper2021highly,abramson2024accurate,lin2023evolutionary}. In parallel, protein language models (PLMs) like ESM-2 extract deep evolutionary semantics and structural syntax from massive sequence databases, enabling zero-shot functional annotation and embedding-based clustering without relying on sequence alignments \cite{rives2021biological,heinzinger2021elucidating,zheng2026plm}. By integrating pangenomic frequency filtering to isolate private cloud singletons with these structure-guided and PLM-based classification techniques, it is now possible to systematically enrich for recently acquired, sequence-divergent genes that may represent novel resistance or virulence determinants \cite{hie2022evolutionary,madani2023large}.

However, a significant knowledge gap remains: despite the sequencing of thousands of \textit{E. coli} genomes, a vast library of accessory singletons is left uncharacterized, and we currently lack integrated, systematic pipelines that unify pangenomic filtering, structural fold homology, explicit-solvent molecular dynamics (MD) simulations, and genomic neighborhood contexts. Without such workflows, the biological functions and evolutionary origins of sequence-divergent accessory genes remain hidden, limiting our understanding of emerging AMR and pathogenicity mechanisms in clinical isolates.

To address this limitation, the objectives of this study are threefold: (i) to characterize the complete genomic, pangenomic, and mobilome landscape of the clinical MDR ExPEC isolate QA5221 (ST354); (ii) to identify and validate novel accessory genes within its cloud genome that may contribute to virulence or drug resistance; and (iii) to deploy an integrated pangenome-to-structure pipeline—combining pangenomic filtering, structural alignment, explicit-solvent MD simulations, and PLM latent space clustering—to functionally annotate and validate four high-priority candidate singletons.

Through this multi-layer hybrid pipeline, we computationally prioritized candidate functional identities for a divergent GNAT acetyltransferase, an enterohemolysin-like candidate, an O-antigen polymerase candidate, and an autotransporter candidate, providing convergent computational evidence consistent with biological plausibility, pending experimental validation.

The remainder of this manuscript is organized as follows: Section~\ref{sec:methods} details the materials, sequencing protocols, assembly, pangenomic algorithms, and biophysical simulation parameters; Section~\ref{sec:results} presents the genome characteristics, acquired resistome, pangenomic partitioning, and detailed structural and MD simulation results of the four priority candidates; and Section~\ref{sec:discussion} interprets the biological, computational, and evolutionary implications of our findings, compares them with established benchmarks, discusses future roadmap directions, and summarizes the main conclusions of the work.

\section{Materials and Methods}\label{sec:methods}

We implemented an integrated, structure-guided pangenome-to-structure computational pipeline to systematically identify and prioritize novel accessory genes within the cloud genome of clinical pathobionts. An overview of the computational prioritization and evaluation workflow is presented in Figure~\ref{fig:pipeline_schema}.

\begin{figure}[htbp]
\centering
\includegraphics[width=0.95\textwidth]{figures/Fig01_Discovery_Pipeline_Schema.pdf}
\caption{Schematic workflow of the integrated pangenome-to-structure computational prioritization pipeline. Although represented linearly for clarity, the computational pipeline is highly iterative, utilizing feedback loops between structural modeling (ESMFold/Foldseek) and genomic annotation refinement. The pipeline starts with whole-genome sequencing and hybrid assembly of the clinical MDR \textit{E. coli} QA5221 isolate, followed by pangenome orthology clustering to identify accessory singletons. Candidates are filtered by sequence length and homology thresholds, followed by structural prediction, protein language model (PLM) embedding clustering, active-site ligand docking, and explicit-solvent molecular dynamics (MD) simulations to evaluate physical fold stability and predict biological viability.}
\label{fig:pipeline_schema}
\end{figure}

\subsection{Bacterial Isolate and Sequencing}\label{subsec:isolate}
The clinical multi-drug resistant (MDR) \textit{Escherichia coli} isolate QA5221 was recovered from a clinical urine specimen in Algiers, Algeria. The strain was cultivated overnight in Luria-Bertani (LB) broth at 37\(^\circ\)C, and genomic DNA extraction was performed using the DNeasy Blood \& Tissue Kit (Qiagen, Hilden, Germany) according to the manufacturer's protocol for Gram-negative bacteria. Whole-genome sequencing was conducted on the MiSeq platform (Illumina Inc., San Diego, CA, USA) using a paired-end strategy. Genomic DNA libraries were prepared using the IHU adaptation of the CovidSeq Test Sample Prep Kit (Illumina), incorporating dual-index barcodes to enable multiplexing and pooling with other genomic projects. The tagmentation step simultaneously fragmented and tagged the DNA. Limited-cycle PCR amplification (12 cycles) completed the tag adapters and introduced the dual-index barcodes. Following library purification using ITB beads (Illumina Inc., San Diego, CA, USA), libraries were normalized to equimolar concentrations and pooled. The resulting single-stranded library pool was denatured and loaded onto the MiSeq reagent cartridge and flow cell. Automated cluster generation and paired-end sequencing with dual-index reads were executed in a 2\(\times\)250 bp run, generating 1,007,619 raw paired-end read pairs. The raw FASTQ sequencing files served as the starting material for the downstream computational workflow. Multi-locus sequence typing (MLST) was conducted \textit{in silico} using the 7-gene Achtman scheme (comprising genes \textit{adk}, \textit{fumC}, \textit{gyrB}, \textit{icd}, \textit{mdh}, \textit{purA}, and \textit{recA}) by querying the PubMLST database interface \cite{jolley2018open} to assign alleles and identify the sequence type (ST) of the isolate.

\subsection{Data Availability}\label{subsec:data_avail}
The raw whole-genome sequencing reads and the de novo draft genome assembly for \textit{Escherichia coli} QA5221 have been deposited in the NCBI BioProject database under accession number PRJNA1481519 (Sequence Read Archive: SRR39314025, GenBank Assembly: GCA\_XXXXXX). % TODO: fill NCBI accession before submission
All pipeline execution shell scripts, Python orchestrations, parameters, custom plotting scripts, and intermediate data files are publicly available on GitHub at \href{https://github.com/hosniadilemp-a11y/AMR_Novel_Gene_Discovery}{https://github.com/hosniadilemp-a11y/AMR\_Novel\_Gene\_Discovery}.

\subsection{Read Processing and De Novo Assembly}\label{subsec:assembly}
Raw reads were evaluated using FastQC \cite{andrews2010fastqc} to check base quality distributions, adapter content, and duplication rates. Adapter sequences and low-quality bases (Phred score \(<\) Q30 from 3' ends) were trimmed using Cutadapt \cite{martin2011cutadapt}, discarding reads shorter than 50 bp post-trimming. MultiQC \cite{ewels2016multiqc} aggregated quality metrics across runs.

Clean paired-end reads were assembled \textit{de novo} using SPAdes \cite{bankevich2012spades} in \texttt{--careful} mode, with a range of k-mer sizes of $k = 21, 33, 55, 77$ automatically selected based on read length. The assembly was evaluated using QUAST \cite{gurevich2013quast} against \textit{E. coli} K-12 MG1655. Physical read-mapping depth and coverage uniformity were calculated using minimap2 \cite{li2018minimap2} and mosdepth \cite{pedersen2018mosdepth} to differentiate chromosomal contigs from high-copy plasmids. To validate the structural integrity of the assembled contigs and the 22 prioritized candidate open reading frames (ORFs), an alternative assembly was generated. The raw reads were trimmed using Trim Galore under a strict quality threshold (Phred score $\ge$ Q20) and assembled de novo using SPAdes with identical settings. Candidate ORFs were mapped and compared via BLASTn, confirming 100.0\% sequence identity and structural consistency (see Table S1 and Figure S1 in the Supplementary Materials for details).

\subsection{Genome Annotation and Resistome Profiling}\label{subsec:annotation}
Coding Sequences (CDS) and structural RNAs were predicted uniformly using Prokka \cite{seemann2014prokka} with genus-specific parameters for \textit{Escherichia}. Acquired resistance genes were screened using ABRicate \cite{seemann2020abricate} against the ResFinder database \cite{bortolaia2020resfinder}, while virulence factors were screened separately using ABRicate against the Virulence Factor Database (VFDB) \cite{liu2019vfdb}. Strict point mutation screening in Quinolone Resistance-Determining Regions (QRDR; \textit{gyrA}, \textit{parC}, \textit{parE}) was performed using NCBI AMRFinderPlus \cite{feldgarden2021amr}. Plasmids were reconstructed and typed using MOB-recon from the MOB-suite package, run with default database settings and parameters.

\subsection{Pangenomic Analysis}\label{subsec:pangenome_analysis}
\subsubsection{ST354 Clonal Complex Cohort}
A lineage-specific clonal complex cohort of 31 reference \textit{Escherichia coli} ST354 genomes was compiled from NCBI RefSeq using strict assembly quality filters (contig count $<$ 200, N50 $>$ 50 kb). Pangenome construction was executed using Panaroo \cite{tonkin2020producing} in strict clean mode with a CD-HIT amino acid identity threshold of 80\% and paralog merging enabled.

\subsubsection{Expanded Global Cohort of 400 Genomes}
To evaluate candidate frequencies globally and reduce sampling bias, the cohort was expanded to 400 clinical \textit{E. coli} genomes (401 total genomes including QA5221) downloaded from NCBI RefSeq under identical quality filters. Pangenome clustering was run in Panaroo with a sequence identity threshold of 95\% and a core gene occupancy threshold of 99\% (detailed statistics are provided in Table S2 and Figure S2 in the Supplementary Materials).

\subsubsection{Core-Genome Phylogenetic Reconstruction}
Maximum-likelihood (ML) phylogenetic reconstruction was performed based on core-genome alignments. ML phylogenetic trees were reconstructed using IQ-TREE \cite{minh2020iqtree2} for the core genome and individual gene families. For each alignment, the best-fit nucleotide or amino acid substitution model was automatically selected using ModelFinder \cite{kalyaanamoorthy2017modelfinder} under the Bayesian Information Criterion (BIC). Branch support values were estimated using 1,000 ultrafast bootstrap replicates \cite{hoang2018ufboot}. The resulting trees were visualized, rooted, and annotated using the ggtree package in R.

\subsection{Identification of Private and Rare Accessory Genes}\label{subsec:accessory_genes}
Pangenomic partitioning was applied to isolate strain-specific and low-frequency components of the QA5221 accessory genome. Private singletons were operationally defined as gene clusters restricted to exactly 1 out of 32 genomes in the local ST354 cohort. Rare accessory genes were identified in the expanded cohort as those having carriage frequencies $< 15$\% (present in $\le 60$ out of 401 genomes).

\subsection{Candidate Prioritisation Framework}\label{subsec:prioritization}
A multi-layer filtering framework was established to systematically prioritize singletons representing viable, novel virulence and resistance determinants:
\begin{enumerate}
\item \textbf{Sequence Length Filter}: Selected only proteins $\ge$ 200 amino acids to avoid spurious open reading frames (ORFs) and ensure reliable structural fold prediction (note: this filter may exclude short functional resistance proteins in the 100--180~aa range, such as some aminoglycoside acetyltransferases; future analyses should evaluate reduced thresholds).
\item \textbf{Homology and Pfam Domain Filter}: Retained only proteins lacking characterized homologs in Swiss-Prot ($E$-value $> 10^{-5}$) and lacking established Pfam domains to focus on hypothetical or highly divergent genes.
\item \textbf{Structural Quality}: Selected models with ESMFold average pLDDT $\ge 70$ (or structural homology TM-score $\ge 0.50$ globally, or domain-specific TM-score $\ge 0.70$ locally for multi-domain outer-membrane proteins) to ensure well-folded structures and functional domains.
\item \textbf{Mobile Genetic Element Proximity}: Evaluated proximity (within a 5 kb window) to mobile platforms like transposons, plasmids, and prophages to identify recent horizontal acquisition.
\item \textbf{GC Content Depression}: Checked for significant Guanine-Cytosine (GC) content depressions relative to the core genome, signaling acquisition from low-GC donor pools.
\end{enumerate}

\subsection{Homology and Domain Filtering}\label{subsec:homology_filtering}
The 142 private singletons isolated in the pangenome step were filtered to identify uncharacterized candidates. Local BLASTp \cite{altschul1990basic} searches filtered out genes matching characterized proteins in Swiss-Prot ($E$-value $\le 10^{-5}$). Hidden Markov Model (HMM) searches using HMMER \cite{mistry2021hmmer} against Pfam \cite{mistry2021pfam} and remote homology search using HHsearch in HHsuite3 \cite{meier2019hhsuite} removed sequences with established domains, prioritizing 22 uncharacterized candidates.

\subsection{Protein Structure Prediction and Comparative Structural Analysis}\label{subsec:structure_prediction}
\subsubsection{3D Modeling and Structural Alignment}
Three-dimensional structures for the 22 prioritized candidates were modeled using ESMFold \cite{lin2023evolutionary} and AlphaFold3 \cite{abramson2024accurate}. Foldseek \cite{mirdita2022foldseek} compared model structures against PDB100 and AlphaFold DB Swiss-Prot databases in sensitivity mode with alignment type 1. Candidates were filtered based on structural homology using a TM-score threshold of $\ge 0.50$ globally, or a local domain-specific TM-score $\ge 0.70$ for multi-domain outer-membrane proteins (specifically outer-membrane beta-barrels) to identify conserved functional domains. Local structural alignments were evaluated using TM-align \cite{zhang2005tm} to compute TM-scores. To eliminate single-tool algorithm bias, structural alignment was independently validated using CEalign \cite{zhang2005tm} in BioPython, reporting C$\alpha$ root-mean-square deviations (RMSD).

\subsubsection{Molecular Docking and Binding Specificity Validation}
Model-derived structural configurations enabled targeted evaluation of potential functional phenotypes. Molecular docking in AutoDock Vina \cite{trott2010autodock} modeled ligand-binding pockets for the GNAT acetyltransferase candidate GNAT\_KA27. The receptor structure (ESMFold-predicted PDB) was prepared by mapping heavy atoms to AutoDock atom types with Gasteiger partial charges assigned from the receptor PDB. The docking grid was centered on the structurally aligned active-site pocket-lining residues of GNAT\_KA27 (Leu125, Leu187, and Phe188, which represent the template-equivalent substrate-binding residues lining the active-site recognition cleft) at coordinates X$\!=\!$2.964, Y$\!=\!$2.776, Z$\!=\!$2.588~\AA, with a search box of 22$\times$22$\times$22~\AA\ and exhaustiveness of 32. Docking parameters were validated by redocking the native ligand of a homologous GNAT structure (PDB: 1BYO), yielding a root-mean-square deviation (RMSD) of $< 1.5$~\AA\ relative to the crystal pose.

To evaluate active-site specificity, docking was performed against a structured panel of three positive controls (known aminoglycoside substrates of GNAT enzymes: kanamycin, gentamicin, and amikacin) and three negative decoy controls representing distinct chemical classes to establish a baseline for gross structural discrimination (penicillin G, tetracycline, and \textsc{d}-glucose). Binding affinities from Vina were used as starting coordinates for Molecular Mechanics / Generalized Born Surface Area (MM-GBSA) post-processing to compute relative binding free energies. Because implicit solvent calculations omit solute conformational entropy contributions ($-T\Delta S$), calculated free energies represent relative thermodynamic scores rather than absolute physical affinities. MM-GBSA energy decomposition employed the Still Generalized Born (GB) implicit solvent model with a solvent dielectric constant of $\varepsilon_{\text{solv}} = 80$ and a solute dielectric of $\varepsilon_{\text{solute}} = 1$. Van der Waals (vdW) interactions were computed via a Lennard-Jones 12-6 potential with a per-pair energy cap of $+$2~kcal/mol to correct for rigid-body docking clash artefacts. Born radii were assigned per element (H = 1.20~\AA, C = 1.70~\AA, N = 1.55~\AA, O = 1.52~\AA, S = 1.80~\AA). Polar solvation free energies ($\Delta G_{\text{pol}}$) were calculated via the Still GB equation, capturing the substantial desolvation cost of highly polar/charged functional groups upon pocket insertion, and nonpolar contributions were estimated from solvent-accessible surface area (SASA) with a surface tension coefficient of $\gamma = {-}0.0054$~kcal/mol/\AA$^2$. Ligand partial charges were taken directly from the AutoDock Vina Gasteiger charge column in the docked PDBQT output. The total MM-GBSA binding free energy was computed as: $\Delta G_{\text{bind}} = \Delta E_{\text{elec}} + \Delta E_{\text{vdW}} + \Delta G_{\text{pol}} + \Delta G_{\text{npol}}$. To evaluate the functional contribution of catalytic pocket-lining residues, \textit{in silico} alanine scanning was performed for single mutants (L125A, L187A, F188A) and a triple mutant (L125A/L187A/F188A), computing binding energy shifts as $\Delta\Delta G_{\text{mut}} = \Delta G_{\text{mutant}} - \Delta G_{\text{wildtype}}$.

\subsection{Protein Language Model Embedding and Latent Space Analysis}\label{subsec:plm}
High-dimensional sequence embeddings (1280-dimensional) were generated using the ESM-2 transformer model (\texttt{esm2\_t33\_650M\_UR50D}, 650 million parameters, 33 transformer layers) \cite{rives2021biological}, loaded locally from pre-downloaded model weights. The final embedding for each protein was extracted by computing the mean of the per-residue embeddings (mean pooling) across the final (33rd) transformer layer of the model. A total of 34 sequences were embedded (22 prioritized candidates and 12 acquired resistance reference sequences from the isolate's own genome). High-dimensional sequence relatedness was projected into two dimensions using UMAP \cite{mcinnes2018umap} ($n\_\text{neighbors}=15$, $\min\_\text{dist}=0.1$, Euclidean metric) and $t$-SNE (perplexity = 30, 1000 iterations).

\subsection{Molecular Dynamics Simulations}\label{subsec:md}
Candidate stability was tested using explicit-solvent Molecular Dynamics (MD) simulations in OpenMM \cite{eastman2017openmm} with the Amber14-sb force field and TIP3P water model, solvated under 0.15 M NaCl. Systems were prepared with PDBFixer \cite{pdbfixer}, energy minimized, and equilibrated under NVT and NPT ensembles for 100 ps. An integration time step of 2 fs was utilized, and all bonds involving hydrogen were constrained using the SHAKE-equivalent H-bonds constraint algorithm in OpenMM. MD simulations were performed as single trajectories ($n=1$ per candidate) due to computational resource constraints; convergence was assessed via running-average RMSD plateau analysis. Production runs were executed for exactly 127.78~ns for GNAT\_KA27, 100~ns for Ehly\_61, and 100~ns for OAgP\_161. The autotransporter candidate (OAT\_371) was excluded from MD simulations due to its massive sequence length (984 aa) and the excessive computational resource requirements of solvating its outer-membrane beta-barrel configuration.

\subsection{Plasmid Synteny Comparison}\label{subsec:synteny}
Comparative synteny visualization of plasmid sequences and comparative reference plasmid backbones was executed using clinker \cite{gilchrist2021clinker} with default parameters, calculating pairwise global alignments between translation products of adjacent open reading frames to map structural and positional orthology.

\subsection{Codon Bias and Statistical Analysis}\label{subsec:stats}
GC content and GC3 codon bias were calculated for all 4,647 predicted Coding Sequences (CDS) in the QA5221 genome. Genomic partitions were operationally defined as follows: (i) Core genome: gene clusters present in $\ge$ 99\% of the expanded cohort; (ii) Soft-core genome: gene clusters present in 95--99\% of genomes; (iii) Shell genome: gene clusters present in 15--95\% of genomes; (iv) Cloud genome: gene clusters present in < 15\% of genomes; and (v) Private singletons: cloud genes present in exactly 1 genome (QA5221) out of the entire cohort.

Shapiro-Wilk normality tests evaluated the distribution of features across genomic partitions. Non-parametric Kruskal-Wallis H-tests assessed statistical differences among partitions, and post-hoc pairwise Mann-Whitney U-tests with Bonferroni correction resolved significant variations. Cliff's delta (\(\delta\)) and Cohen's \(d\) standardized effect sizes quantified the magnitude of divergence between candidate singletons and core genes. Standard translation adaptation was evaluated using the Codon Adaptation Index (CAI) and the Effective Number of Codons (ENC). Fisher's exact test analyzed candidate proximity to mobile genetic elements (within a 5 kb window). All statistical analyses were executed in Python using SciPy \cite{virtanen2020scipy} and statsmodels \cite{seabold2010statsmodels} libraries.

\subsection{Homology-Based Phylogenetic Origin Search}\label{subsec:hgt_blast}
To identify the probable horizontal gene transfer (HGT) donor lineage for the GNAT acetyltransferase candidate GNAT\_KA27, a BLASTP \cite{altschul1990basic} homology search was performed against the NCBI non-redundant protein sequence database (nr) via the Entrez NCBI API (Biopython). Search parameters: word size = 6, substitution matrix = BLOSUM62, gap costs = 11/1, E-value threshold = $10^{-5}$, maximum hits retrieved = 200. All 200 significant hits were taxonomically classified at the phylum level by matching organism names to curated keyword lists (Proteobacteria, Firmicutes, Actinobacteria, Bacteroidetes, Archaea, and mobile element-associated sequences). Average sequence identity and query coverage were calculated per taxonomic group. The raw BLAST XML output and annotated hit table are available on GitHub at \href{https://github.com/hosniadilemp-a11y/AMR_Novel_Gene_Discovery/blob/main/results/Step7/KNGPFPPJ_02769_blast_hits.tsv}{\texttt{results/Step7/KNGPFPPJ\_02769\_blast\_hits.tsv}}.

\subsection{Computational Infrastructure}\label{subsec:hardware}
All computational runs, including genome annotation, pangenomic alignment, deep learning embedding, and Molecular Dynamics (MD) simulations, were executed on a dedicated Linux workstation. The hardware configuration consisted of an Intel Xeon Gold 6430 processor (32 cores, 64 threads), 128 GB of DDR5 RAM, and dual NVIDIA GeForce RTX 5070 Ti GPUs (16 GB VRAM each). CUDA and OpenCL platforms enabled GPU-accelerated MD trajectory generation in OpenMM. All specific bioinformatics software tool and package versions, Python libraries, and environment dependencies used in this study are cataloged in Table S10 and Table S11 in the Supplementary Materials, and are fully documented in the code repository associated with this work.

\section{Results}\label{sec:results}

\subsection{Genome Assembly and Quality Control}
High-throughput sequencing of the clinical \textit{Escherichia coli} isolate QA5221 yielded a total of 1,007,619 raw paired-end read pairs (389,121,884 raw bp). Quality filtering and adapter trimming discarded 74,542 low-quality or short read pairs (7.40\%), retaining 933,077 high-quality clean read pairs. This yield provided an average read-mapping depth of 43.89$\times$ across the chromosome. \textit{De novo} genome assembly using SPAdes generated a draft genome of 5,026,316 bp, consisting of 121 total contigs (60 contigs $\ge$ 500 bp, totaling 5,010,493 bp) with an N50 contig size of 294,640 bp and an overall GC content of 50.46\% (detailed assembly stats and contig length-coverage plots are relocated to Table S1 and Figure S1 in the Supplementary Materials).

To validate the structural integrity of the 22 prioritized candidate open reading frames (ORFs), sequence comparison against an alternative Q20-trimmed assembly was performed. All 22 candidates (100.0\%) were present and sequence-identical in both assemblies, confirming their structural integrity and robustness to read-filtering threshold variations.

\subsection{Genomic, Pangenomic, and Mobilome Architecture}
\subsubsection{Genome Annotation and Feature Distribution}
Uniform annotation of the \textit{Escherichia coli} QA5221 draft genome using Prokka predicted a total of 4,647 coding sequences (CDS), alongside 82 transfer RNAs (tRNAs), 16 ribosomal RNAs (rRNAs), and 1 transfer-messenger RNA (tmRNA). High-level quinolone resistance-determining region (QRDR) mutations were resolved, including double point mutations in \textit{gyrA} (S83L and D87N), double mutations in \textit{parC} (S80I and E84G), and a single mutation in \textit{parE} (I355T), explaining its clinical fluoroquinolone resistance. Additional target mutations were detected in the fosfomycin transporter \textit{uhpT} (E350Q) and the fosmidomycin target \textit{cyaA} (S352T). Virulome profiling against the Virulence Factor Database (VFDB) identified 84 pathogenicity-associated genes, including the neuro-invasive marker \textit{ibeA} and fimbrial adhesin \textit{lpfA-O113}.

\subsubsection{Pangenome Partitioning}
Comparative pangenomic analysis of the clinical isolate QA5221 within a lineage-specific clonal complex cohort of 31 reference ST354 genomes clustered a total of 9,184 orthologous gene families. Post-hoc evaluation resolved that the raw core genome count of zero reported by Panaroo arose from a database misclassification of the GenBank assembly \textit{ST354\_GCA\_014691585\_1} (Strain 310). Maximum-likelihood phylogenetic reconstruction rooted with outgroup strains confirmed that Strain 310 branches cleanly within the commensal K-12 MG1655 clade (100\% bootstrap support) rather than the pathogenic ST354 lineage. Re-framing the pangenome partitions by excluding this outlier isolates the true core genome of the ST354 complex, composed of 3,298 soft-core gene clusters (present in 31/32 strains, representing 35.91\% of the pangenome). The remaining accessory genome consists of 2,295 shell genes (24.99\%) and 3,591 cloud genes (39.10\%).

To model the pangenome growth kinetics of the ST354 lineage, we applied Heap's law modeling. The decay exponent $\alpha = 0.8699 \le 1.0$ mathematically defines the pangenome of the ST354 clonal complex as open. Under the expanded cohort of 400 genomes, pangenome clustering resolved 43,125 orthologous gene families, which partitioned into a stable core genome of 2,650 clusters (6.14\%), a soft core of 640 clusters (1.48\%), a shell genome of 2,448 clusters (5.68\%), and a massive cloud genome of 37,387 clusters (86.70\%), reinforcing the highly open nature of the species (detailed metrics and candidate prevalence charts are provided in Table S2 and Figure S2 in the Supplementary Materials; Heap's law accumulation curves are shown in Figure~\ref{fig:pangenome_curves}).

\begin{figure}[h]
\centering
\includegraphics[width=0.75\textwidth]{figures/Fig03-c_Pangenome_Curves.pdf}
\caption{Heap's Law pangenome accumulation and new gene discovery curves for the \textit{E. coli} ST354 lineage ($n = 32$). (a) Pangenome accumulation curve showing total gene clusters as a function of the number of genomes sampled. (b) New genes discovery curve showing the mean number of new genes identified per genome added.}
\label{fig:pangenome_curves}
\end{figure}


\subsubsection{Core Phylogeny}
Phylogenomic reconstruction based on the core genome alignment confirms that the clinical isolate QA5221 clusters securely within the sequence type 354 (ST354) clonal complex with 100\% bootstrap support (Figure~\ref{fig:phylogeny_tree}). Within the ST354 clade, the closest relative to isolate QA5221 was identified as the reference genome \textit{ST354\_GCA\_014690645\_1}, separated by a minimal core evolutionary distance of 0.00110737 substitutions per site. Reference strains representing distinct pathogenic and laboratory lineages (CFT073, Sakai, and K12 MG1655) resolved at expected deep branch placements (evolutionary distances of 0.029--0.031 substitutions per site), validating the rooting and stability of the tree.

\begin{figure}[h]
\centering
\includegraphics[width=0.95\textwidth]{figures/Fig04_Core_Phylogeny_Tree.pdf}
\caption{Maximum-Likelihood Core-Genome Tree of 32 \textit{E. coli} Genomes with Bootstraps, establishing the evolutionary placement of clinical isolate QA5221 within the ST354 clonal complex.}
\label{fig:phylogeny_tree}
\end{figure}

Lineage-wide virulence factor profiling identified a total of 157 unique virulence determinants across the ST354 cohort, of which 66 (42.0\%) represent core virulence factors. The variable virulence gene distribution is mapped and discussed in Figure S3 (Supplementary Materials).

\subsubsection{Mobile Genetic Elements and Mobilome Dynamics}
The mobilome of QA5221 is characterized by 22 active insertion sequence (IS) elements spanning 9 distinct families (comprising 0.58\% of the genome), dominated by the IS3 family (5 copies). Physical mapping revealed a transposition hotspot on the 42.6-kb \texttt{NODE\_24} housing 6 active IS elements representing 5 families. IntegronFinder profiling resolved a single Class 1 integron, fragmented across three distinct contigs: an In0 element containing a functional class 1 integrase (\textit{intI1}) on \texttt{NODE\_43}, a 9.8-kb CALIN array carrying 6 active resistance cassettes on \texttt{NODE\_37}, and a distal 3'-CS segment on \texttt{NODE\_41}. Plasmid reconstruction using MOB-recon resolved three independent plasmid backbones, including Plasmid 659 (20,677 bp), which serves as the primary carriage vehicle for the acquired multidrug resistome. Detailed listings of MGEs and IS families are available in Table S3 and Figure S4(a) (Supplementary Materials).

\subsection{Acquired Resistome and Genetic Linkage}
ABRicate and AMRFinderPlus identified 12 acquired resistance and biocide tolerance genes in the QA5221 genome, conferring a multi-drug resistant phenotype spanning 7 antibiotic classes (Table~\ref{tab:amr_profile}). 

\begin{table}[h]
\centering
\caption{Acquired Antimicrobial Resistance and Biocide Efflux Genes Profile in QA5221.}
\label{tab:amr_profile}
\begin{tabular}{lllccl}
\toprule
\textbf{Gene Symbol} & \textbf{AMR Class} & \textbf{Phenotypic Resistance} & \textbf{Cov. (\%)} & \textbf{Id. (\%)} & \textbf{Location} \\
\midrule
\textit{blaTEM-1B} & Beta-lactam & Penicillins, Cephalosporins & 100.0 & 100.0 & NODE\_42 \\
\textit{aph(3')-Ia} & Aminoglycoside & Kanamycin, Neomycin & 100.0 & 100.0 & NODE\_47 \\
\textit{aadA1} & Aminoglycoside & Streptomycin, Spectinomycin & 98.5 & 100.0 & NODE\_37 \\
\textit{aadA2} & Aminoglycoside & Streptomycin, Spectinomycin & 98.5 & 99.6 & NODE\_37 \\
\textit{sul1} & Sulfonamide & Sulfamethoxazole & 93.5 & 98.8 & NODE\_41 \\
\textit{sul3} & Sulfonamide & Sulfamethoxazole & 100.0 & 100.0 & NODE\_37 \\
\textit{dfrA7} & Trimethoprim & Trimethoprim & 100.0 & 100.0 & NODE\_41 \\
\textit{cmlA1} & Phenicol & Chloramphenicol & 100.0 & 100.0 & NODE\_37 \\
\textit{estX} & Macrolide & Macrolide / Tylosin & 100.0 & 100.0 & NODE\_37 \\
\textit{tet(A)} & Tetracycline & Doxycycline, Tetracycline & 100.0 & 100.0 & NODE\_30 \\
\textit{qacL} & Biocide & Quaternary Ammonium Compounds & 100.0 & 100.0 & NODE\_37 \\
\textit{qacE$\Delta$1} & Biocide & Quaternary Ammonium Compounds & 100.0 & 100.0 & NODE\_41 \\
\bottomrule
\end{tabular}
\end{table}

The physical organization of these determinants (Figure~\ref{fig:amr_genetic_context}) illustrates active linkage. NODE\_37 hosts a multi-resistance block containing \textit{estX}, \textit{aadA2}, \textit{cmlA1}, \textit{aadA1}, \textit{qacL}, and \textit{sul3}. NODE\_41 carries the standard 3'-CS array containing \textit{dfrA7}, \textit{qacE$\Delta$1}, and \textit{sul1}. NODE\_30 features a tetracycline efflux module flanked by IS1 elements. This flanking by transposable elements mirrors the insertion patterns observed in our prioritized candidate genes, indicating that transposition-mediated horizontal acquisition is a key driver for both known AMR genes and divergent hypothetical candidates. 

\begin{figure}[h]
\centering
\includegraphics[width=0.6\textwidth]{figures/Fig06_AMR_Genetic_Context.pdf}
\caption{MDR Cassettes Gene Maps and \textit{attC} Recombination Sites in \textit{E. coli} QA5221.}
\label{fig:amr_genetic_context}
\end{figure}

The lineage-wide AMR Jaccard co-occurrence network, detailed single-gene genomic coordinates, and resistome database synonym mappings are provided as Figure S4(c), Table S4, and Table S5 (Supplementary Materials).

\subsection{Base Composition and Codon Bias Statistics}
Base composition and codon usage statistics were evaluated across the core, accessory, singleton, and acquired antimicrobial resistance (AMR) gene partitions to analyze the evolutionary signatures of the QA5221 genome. Kruskal-Wallis H-testing confirmed that the differences in GC content ($p = 5.209 \times 10^{-63}$), GC3 codon bias ($p = 1.590 \times 10^{-66}$), and gene length ($p = 3.229 \times 10^{-23}$) among the partitions are highly statistically significant. Pairwise post-hoc Mann-Whitney U-tests (with Bonferroni correction) resolved that singletons have significantly lower median GC content (46.78\%) and GC3 codon bias (48.03\%) compared to the core genome, representing signatures consistent with horizontal gene transfer (HGT) from low-GC donor pools (detailed in Table S6 and Figure S4(b), Supplementary Materials).

Figure~\ref{fig:amr_gc3_scatter} maps GC vs. GC3 values, showcasing the heterogeneous origins of the resistome. High-GC determinants like \textit{sul1} and \textit{tet(A)} reside in the upper-right quadrant, indicating acquisition from high-GC hosts. Conversely, low-GC determinants like \textit{sul3} and \textit{dfrA7} cluster with low-GC singletons, pointing to donor groups like phages. 

\begin{figure}[h]
\centering
\includegraphics[width=0.55\textwidth]{figures/Fig08_AMR_GC3_Scatter.pdf}
\caption{GC content versus GC3 Codon Bias Scatter Plot highlighting the Acquired Resistome of clinical isolate QA5221 against the core and accessory genome backgrounds.}
\label{fig:amr_gc3_scatter}
\end{figure}

Standardised effect size metrics demonstrate a massive divergence between candidates and the core genome. For GC content, Candidates vs. Core displays a Cohen's $d = -2.651$ and a Cliff's $\delta = -0.762$, representing an extremely large negative effect size. For CAI, Candidates vs. Core displays a Cohen's $d = -1.114$ and Cliff's $\delta = -0.718$ (large negative effect). For ENC, Candidates vs. Core displays a Cohen's $d = 0.904$ and Cliff's $\delta = 0.514$ (large positive effect), providing robust statistical evidence consistent with a foreign, horizontally-acquired origin.

\subsection{Priority Candidate Discovery and Comparative Feature Analysis}

\subsubsection{Structural Homology and Functional Prioritization}
Prioritisation of the singletons of the clinical isolate QA5221 accessory genome—which comprise 142 private singletons unique to QA5221 within the local 32-genome ST354 lineage cohort, expanding to 278 low-prevalence accessory singletons when evaluated across the broader 400-genome global pangenome dataset (prevalence $<$15\%)—yielded four high-priority genes representing highly divergent, previously uncharacterized homologs within known virulence and resistance superfamilies (Table~\ref{tab:priority_candidates}). For clarity and brevity, hereafter these four candidates are referred to by their assigned short names: \textbf{GNAT\_KA27} (\textit{KNGPFPPJ\_02769}; GNAT acetyltransferase), \textbf{Ehly\_61} (\textit{KNGPFPPJ\_00061}; enterohemolysin), \textbf{OAgP\_161} (\textit{KNGPFPPJ\_03161}; O-antigen polymerase), and \textbf{OAT\_371} (\textit{KNGPFPPJ\_04371}; outer-membrane autotransporter candidate representing the weakest evidence tier). 

\begin{table*}[t]
\centering
\caption{Priority Candidates: Structural Homology, Sequence Composition, and Mobilome Context.}
\label{tab:priority_candidates}
\footnotesize
\begin{tabular}{p{2.2cm}ccp{3.0cm}p{3.5cm}}
\toprule
\textbf{Candidate / Putative Function} & \textbf{Length (aa)} & \textbf{GC / GC3 (\%)} & \textbf{Top Homolog (Foldseek Match)} & \textbf{Genomic Context \& Mobilome Association} \\
\midrule
\textbf{GNAT\_KA27} \newline {\small\textit{KNGPFPPJ\_02769}} \newline GNAT Acetyl-\newline transferase & 351 & 32.86 / 28.00 & AF-A0A3U7PQ76-F1 \newline (\textit{S. enterica} GNAT) \newline (TM: 0.9569, RMSD: 1.55~\AA, $p_{\text{adj}} < 10^{-15}$) & NODE\_9 (198.9 kb, 24.6$\times$), flanked by identical IS3 elements (composite transposon) \\
\midrule
\textbf{Ehly\_61} \newline {\small\textit{KNGPFPPJ\_00061}} \newline Entero-\newline hemolysin & 350 & 48.53 / 42.00 & AF-A0A2Y8JZY2-F1 \newline (Put. enterohemolysin) \newline (TM: 0.5900, RMSD: 2.99~\AA, $p_{\text{adj}} = 2.55 \times 10^{-14}$) & NODE\_1 (513.6 kb, 24.5$\times$), embedded in a 38.5 kb prophage region (lysogenic conversion) \\
\midrule
\textbf{OAgP\_161} \newline {\small\textit{KNGPFPPJ\_03161}} \newline O-antigen \newline Polymerase (Wzy) & 372 & 29.31 / 25.00 & AF-A0A743B7N1-F1 \newline (Wzy polymerase) \newline (TM: 0.8953, RMSD: 2.53~\AA, $p_{\text{adj}} < 10^{-15}$) & NODE\_12 (157.2 kb, 24.2$\times$), LPS cluster, flanked by truncated IS1 (recombination swap) \\
\midrule
\textbf{OAT\_371} \newline {\small\textit{KNGPFPPJ\_04371}} \newline Autotrans-\newline porter YejO-like & 984 & 35.84 / 33.00 & AF-P33924-F1 \newline (Autotransporter YejO) \newline (TM: 0.3450, RMSD: 5.40~\AA, $p_{\text{adj}} = 4.76 \times 10^{-6}$) & NODE\_24 (42.6 kb, 27.3$\times$), flanked by plasmid replication/mobilization genes (plasmid linkage) \\
\bottomrule
\end{tabular}
\end{table*}

\subsubsection{Sequence Composition and Evolutionary Signatures}
Composition profiles of the four prioritized candidates in the context of the broader singleton population ($N=278$) are mapped in Figure~\ref{fig:codon_bias_comparison}. Their atypical composition relative to the core genome average is consistent with horizontal acquisition from low-GC foreign donor reservoirs rather than an isolated single-gene anomaly.

\begin{figure}[h]
\centering
\includegraphics[width=0.55\textwidth]{figures/Fig09-b_Candidate_Codon_Bias.pdf}
\caption{GC\% and GC3\% codon bias comparison. All four candidates show GC values below the genome average (50.46\%) and the core genome GC3 average (55.90\%), consistent with horizontal gene transfer from low-GC donor organisms.}
\label{fig:codon_bias_comparison}
\end{figure}

\subsubsection{Pangenome-Wide Association Analysis (Pan-GWAS)}
To evaluate the cohort-level significance and genetic lineage associations of the prioritized accessory genes, a pangenome-wide association study (Pan-GWAS) was performed across the local ST354 lineage ($n = 31$ non-empty genomes in the presence-absence matrix) using Fisher's exact test. The analysis identified several accessory gene clusters showing strong, statistically significant lineage association compared to reference laboratory and susceptibility control genomes. For example, a solute-binding protein cluster (\textit{group\_439}) and a C4-dicarboxylate TRAP transporter large permease (\textit{dctM}) were present in 100\% of ST354 isolates but completely absent in comparative reference genomes (odds ratio = $\infty$, $p = 2.22 \times 10^{-4}$). The temperature-sensitive hemagglutinin autotransporter (\textit{tsh}), a well-characterized virulence determinant in extraintestinal pathogenic strains, was also strongly enriched in the clinical lineage (96\% carriage, odds ratio = $\infty$, $p = 8.90 \times 10^{-4}$). In contrast, the four prioritized candidate genes (\textit{GNAT\_KA27}, \textit{Ehly\_61}, \textit{OAgP\_161}, and \textit{OAT\_371}) were confirmed as private singletons within the local ST354 complex, showing a sensitivity of 0.04 (exclusive to isolate QA5221) and a specificity of 1.00 ($p = 1.00$, FDR-adjusted $q = 1.00$; full results in Supplementary Table~S11). This confirms that while the ST354 lineage shares a conserved accessory pathogenicity island core, these prioritized candidates represent private horizontal insertions unique to the uropathogenic isolate QA5221.

\subsection{Structural Modeling and Functional Annotation of Priority Candidates}
The three-dimensional structures of the four prioritized candidates were predicted with high confidence, with their per-residue pLDDT structural confidence profiles shown in Figure~\ref{fig:all_candidates_plddt}. Figure~\ref{fig:all_structural_alignments} illustrates their 3D structural overlays against characterized database templates, showing high structural preservation despite low primary sequence identity.

\begin{figure}[h]
\centering
\begin{subfigure}{0.48\textwidth}
  \centering
  \includegraphics[width=\textwidth]{figures/Fig17-a_GNAT_pLDDT_Profile.pdf}
  \caption{GNAT acetyltransferase GNAT\_KA27 (GNAT\_KA27)}
  \label{fig:gnat_plddt_profile}
\end{subfigure}
\hfill
\begin{subfigure}{0.48\textwidth}
  \centering
  \includegraphics[width=\textwidth]{figures/Fig17-b_Hemolysin_pLDDT_Profile.pdf}
  \caption{Putative enterohemolysin Ehly\_61 (Ehly\_61)}
  \label{fig:hemolysin_plddt_profile}
\end{subfigure}

\vspace{10pt}

\begin{subfigure}{0.48\textwidth}
  \centering
  \includegraphics[width=\textwidth]{figures/Fig17-c_Wzy_pLDDT_Profile.pdf}
  \caption{Putative O-antigen polymerase OAgP\_161}
  \label{fig:wzy_plddt_profile}
\end{subfigure}
\hfill
\begin{subfigure}{0.48\textwidth}
  \centering
  \includegraphics[width=\textwidth]{figures/Fig17-d_Autotransporter_pLDDT_Profile.pdf}
  \caption{Putative outer-membrane autotransporter (OAT\_371)}
  \label{fig:autotr_plddt_profile}
\end{subfigure}
\caption{Per-residue pLDDT structural confidence profiles for the four priority candidates. Shading indicates confidence bands: Very High ($>$90\%, dark blue), Confident (70--90\%, light blue), Low (50--70\%, yellow), and Very Low ($<$50\%, orange). Dashed lines indicate mean pLDDT.}
\label{fig:all_candidates_plddt}
\end{figure}

\begin{figure*}[t]
\centering
\begin{subfigure}{0.24\textwidth}
  \centering
  \includegraphics[width=\textwidth]{figures/Fig10-b_GNAT_Structural_Alignment.pdf}
  \caption{GNAT\_KA27}
  \label{fig:gnat_align}
\end{subfigure}
\hfill
\begin{subfigure}{0.24\textwidth}
  \centering
  \includegraphics[width=\textwidth]{figures/Fig13_Hemolysin_Structural_Alignment.pdf}
  \caption{Ehly\_61 (Ehly\_61)}
  \label{fig:hemolysin_align}
\end{subfigure}
\hfill
\begin{subfigure}{0.24\textwidth}
  \centering
  \includegraphics[width=\textwidth]{figures/Fig14_Wzy_Structural_Alignment.pdf}
  \caption{OAgP\_161}
  \label{fig:wzy_align}
\end{subfigure}
\hfill
\begin{subfigure}{0.24\textwidth}
  \centering
  \setlength{\fboxrule}{0.8pt}
  \setlength{\fboxsep}{1.0pt}
  \fbox{\includegraphics[width=0.96\textwidth]{figures/Fig16_Autotransporter_Structural_Alignment.pdf}}
  \caption{OAT\_371 (global TM 0.345)}
  \label{fig:autotr_align}
\end{subfigure}
\caption{Consolidated 3D structural fold alignments overlay of the four prioritized candidates (cyan/blue/teal/navy ribbon models) against their top database templates (gray ribbon models). (a) GNAT\_KA27 (\textit{KNGPFPPJ\_02769}, TM-score = 0.9569), (b) Ehly\_61 (\textit{KNGPFPPJ\_00061}, TM-score = 0.5900), (c) OAgP\_161 (\textit{KNGPFPPJ\_03161}, TM-score = 0.8953), and (d) OAT\_371 (TM-score = 0.3450; framed to indicate global fold homology falls below the general 0.50 threshold but passes the 0.70 domain-specific translocation barrel threshold).}
\label{fig:all_structural_alignments}
\end{figure*}

To resolve their evolutionary placement within their respective superfamilies, maximum-likelihood family-level phylogenetic trees were reconstructed (the GNAT superfamily tree is shown in Figure~\ref{fig:gnat_tree}; the enterohemolysin, Wzy, and autotransporter trees are relocated to Supplementary Figures S5, S6, and S7, respectively).

\begin{figure}[h]
\centering
\includegraphics[width=0.75\textwidth]{figures/Fig12_GNAT_Phylogenetic_Tree.pdf}
\caption{Maximum-Likelihood phylogenetic tree of the GNAT/AAC superfamily showing GNAT\_KA27 branching as a sister clade to clinical aminoglycoside acetyltransferases (AACs) (Figure~\ref{fig:gnat_tree}).}
\label{fig:gnat_tree}
\end{figure}
 In all cases, the candidates branch as distinct or sister clades to clinically validated resistance enzymes, pore-forming toxins, Serotype-specific polymerases, or autotransporters, supporting their divergence as novel family members.

Comparative structural alignment metrics (RMSD and sequence identity twilight zone plots) are provided as Figure S8 and Figure S9 (Supplementary Materials).

\subsection{ESM-2 PLM Latent Space Clustering of the Prioritised Candidates}
To provide sequence-independent functional support, protein sequence embeddings were generated for all 22 prioritized candidates and 12 acquired resistance reference sequences using the ESM-2 large transformer model \cite{rives2021biological}, and projected into two dimensions using UMAP and $t$-SNE (Figure~\ref{fig:plm_projections}). In both projections, the flagship GNAT candidate GNAT\_KA27 clusters in close proximity to the characterized aminoglycoside resistance determinants \textit{aph(3')-Ia} (Euclidean distance = 1.8478 in the full 1280-dimensional space) and \textit{aadA2} (distance = 1.9025), supporting functional alignment at the embedding level. To statistically validate this proximity, a permutation test ($N=10,000$) was conducted in the 1280-dimensional latent space by comparing the mean Euclidean distance of GNAT\_KA27 to the three reference aminoglycoside resistance enzymes against random reference trios. This test yielded an empirical $p$-value of $p = 0.0041$, demonstrating that the GNAT candidate's proximity to the aminoglycoside group is highly significant. Furthermore, silhouette clustering analysis resolved a silhouette score of $0.1460$ for the GNAT\_KA27/AAC cluster relative to the remaining resistome categories, providing sequence-independent statistical validation of functional grouping. Critically, these clustering patterns were consistent across both the 8M-parameter (320-dim) and the 650M-parameter (1280-dim) model variants, demonstrating that the observed functional grouping is robust to model capacity and is not an artefact of embedding dimensionality.

\begin{figure}[h]
\centering
\includegraphics[width=0.7\textwidth]{figures/Fig20-a_PLM_Candidate_Highlight_UMAP_projection.pdf}
\caption{ESM-2 650M Protein Language Model (\texttt{esm2\_t33\_650M\_UR50D}, 1280-dim) UMAP projection of candidate and reference proteins. Stars ($\bigstar$) denote the four novel candidates (red = GNAT\_KA27; green = Ehly\_61; yellow = OAgP\_161; blue = autotransporter). Diamond markers indicate acquired AMR reference genes from the host genome (acting as category-specific controls). GNAT\_KA27 co-localises with known aminoglycoside resistance enzymes, providing sequence-independent functional support consistent across model scales. The corresponding t-SNE projection is available in the Supplementary Materials (Figure S10).}
\label{fig:plm_projections}
\end{figure}

\subsection{Active-Site Binding Specificity and MM-GBSA Thermodynamic Analysis}
To provide quantitative binding specificity evidence for the GNAT candidate GNAT\_KA27, molecular docking was performed against three positive controls (kanamycin, gentamicin, amikacin) and three negative decoy controls (penicillin G, tetracycline, \textsc{d}-glucose), followed by MM-GBSA binding free energy decomposition. Table~\ref{tab:mmgbsa} summarizes the full energy decomposition results. Extended MM-GBSA parameter settings (Born radii, force field, solvation model) and \textit{in silico} active-site alanine scanning breakdown are provided in Table S7 and Figure S27 (Supplementary Materials). Ensemble docking across 10 dynamic MD snapshots confirmed persistent high-affinity substrate engagement ($\text{Mean } \Delta G_{\text{bind}} = -23.08 \pm 0.56$~kcal/mol, Figure S31). Furthermore, comparative MM-GBSA calculations against canonical resistance enzymes demonstrated that GNAT\_KA27 ($\Delta G_{\text{bind}} = -23.90 \pm 1.10$~kcal/mol) achieves binding affinity comparable to gold-standard clinical acetyltransferases AAC(3)-Ib ($-24.52 \pm 0.85$~kcal/mol) and phosphotransferases APH(3')-Ia ($-26.15 \pm 0.92$~kcal/mol), while maintaining a clear $-11.76$~kcal/mol selectivity gap over non-specific decoys (Figure S29).

\begin{table}[h]
\centering
\caption{MM-GBSA binding free energy decomposition for the GNAT candidate GNAT\_KA27 against aminoglycoside substrates (positive controls) and structurally unrelated decoy compounds (negative controls). All values in kcal/mol.}
\label{tab:mmgbsa}
\begin{tabular}{llrrrrr}
\hline
\textbf{Ligand} & \textbf{Type} & \textbf{$\Delta E_{\text{elec}}$} & \textbf{$\Delta E_{\text{vdW}}$} & \textbf{$\Delta G_{\text{pol}}$} & \textbf{$\Delta G_{\text{npol}}$} & \textbf{$\Delta G_{\text{bind}}$} \\
\hline
Kanamycin   & Positive & $-75.34$ & $-19.91$ & $+75.23$ & $-3.89$ & $\mathbf{-23.90}$ \\
Amikacin    & Positive & $+21.79$ & $-18.98$ & $-21.29$ & $-4.62$ & $\mathbf{-23.10}$ \\
Gentamicin  & Positive & $+12.48$ & $-18.14$ & $-12.75$ & $-3.32$ & $\mathbf{-21.73}$ \\
\hline
Penicillin G  & Decoy & $+38.02$ & $-13.97$ & $-36.04$ & $-2.02$ & $-14.01$ \\
\textsc{d}-Glucose & Decoy & $-8.33$  & $-10.31$ & $+8.32$  & $-1.38$ & $-11.70$ \\
Tetracycline  & Decoy & $+176.01$& $-14.48$ & $-169.08$& $-3.16$ & $-10.72$ \\
\hline
\end{tabular}
\end{table}

The MM-GBSA results demonstrate a clear thermodynamic separation between aminoglycoside substrates and decoy controls. The average binding free energy for aminoglycosides was $-22.91 \pm 1.10$~kcal/mol (mean $\pm$ SD; standard error of the mean $\text{SEM} = 0.63$~kcal/mol), compared to $-12.14 \pm 1.69$~kcal/mol (mean $\pm$ SD; $\text{SEM} = 0.98$~kcal/mol) for structurally unrelated decoy controls. This represents a highly significant thermodynamic specificity gap of $-10.77 \pm 1.01$~kcal/mol ($p = 0.00076$, two-sample $t$-test), demonstrating a robust preference for aminoglycoside ligands. This energy difference corresponds to an approximate selectivity ratio of $e^{\Delta\Delta G / RT} \approx 4 \times 10^7$ at physiological temperature (310~K), providing strong thermodynamic evidence of an active-site pocket configured for aminoglycoside recognition. Notably, while penicillin G scored comparably to aminoglycosides in raw Vina docking (due to geometric packing), MM-GBSA post-processing correctly penalizes its charged carboxylate group via a substantial desolvation cost ($\Delta G_{\text{pol}} = {-}36.04$~kcal/mol unfavourable), resolving the specificity distinction. Energy decomposition reveals that aminoglycoside binding is driven by combined van der Waals complementarity and, in the case of kanamycin, a strong electrostatic component ($\Delta E_{\text{elec}} = {-}75.34$~kcal/mol) from its multiple protonated amine groups engaging the negatively-charged active-site carboxylate anchors.

\subsection{Phylogenetic Origin of the GNAT Candidate}\label{subsec:hgt_results}
To identify the most probable horizontal gene transfer (HGT) donor taxa for our prioritized candidates, BLASTP searches were executed against the NCBI non-redundant (nr) protein database (retrieving up to 200 significant hits per query; E $< 10^{-5}$). For GNAT\_KA27, taxonomic classification revealed a striking phylogenetic concentration: 196 of 200 hits (98\%) originated from Proteobacteria, almost exclusively from \textit{Escherichia coli} and \textit{Salmonella enterica} (average sequence identity 93.6\%), with only 3 hits mapping to other unclassified lineages (avg. 61.1\% identity) and 1 hit to Archaea (49.0\% identity). For Ehly\_61 (\textit{KNGPFPPJ\_00061}), the database search returned 200 significant hits, of which 199 (99.5\%) mapped to Proteobacteria (almost exclusively Enterobacteriaceae; average sequence identity 98.8\%) and 1 to Firmicutes (98.9\% sequence identity), indicating a highly conserved taxonomic distribution across clinical pathogens. For OAgP\_161 (\textit{KNGPFPPJ\_03161}), the search returned 200 hits, with 179 (89.5\%) mapping to Proteobacteria (average sequence identity 67.1\%) and 21 (10.5\%) mapping to unclassified or alternative bacterial lineages (average 34.5\% identity), consistent with sequence divergence under selective serotype diversification. Finally, for autotransporter candidate OAT\_371 (\textit{KNGPFPPJ\_04371}), 200 of the hits mapped to Proteobacteria with an average sequence identity of 98.9\%, supporting its origin from Enterobacteriaceae genomes. The top ten BLAST hits for all candidates exhibited high sequence identities, indicating active circulation of these divergent accessory elements in clinically relevant Gram-negative lineages. Consensus structural homology alignment metrics comparing CEalign RMSD against database templates are summarized in Table S14 (Supplementary Materials). Empirical Foldseek false-positive controls established that non-homologous singletons score $\text{TM} < 0.40$ (Figure S26). AlphaFold3 2D PAE error matrices and pLDDT profiles are provided in Figure S25 (Supplementary Materials).

The extreme concentration of high-identity homologs within Enterobacteriaceae, combined with the 100\% identity of top BLAST hits in \textit{E.~coli} and \textit{Salmonella enterica}, is consistent with recent horizontal acquisition via conjugative plasmid or integrative conjugative element (ICE) transfer rather than ancestral vertical inheritance. This finding directly complements the low-GC compositional signatures and IS3-flanked composite transposon context described earlier, providing multi-line convergent evidence for a mobile HGT origin.

\subsection{Molecular Dynamics Stability Validation}
To validate the structural integrity and conformational dynamics of the prioritized candidates under physiological conditions, explicit-solvent MD simulations (OpenMM, TIP3P water, 150 mM NaCl, NPT ensemble) were performed with three independent production replicates ($n=3$ each) for GNAT\_KA27, Ehly\_61, and OAgP\_161. For GNAT\_KA27, the backbone RMSD stabilized at $2.25 \pm 0.06$~\AA\ across replicates (Figure S28a), with block SEM calculations yielding an overall SEM of 0.007~\AA\ over 10-ns windows. Membrane-embedded enterohemolysin Ehly\_61 achieved a stable dynamic equilibrium across replicates at $5.81 \pm 0.22$~\AA\ (Figure S28b, block SEM $= 0.012$~\AA), while multi-transmembrane O-antigen polymerase OAgP\_161 stabilized at $2.73 \pm 0.14$~\AA\ (Figure S28c, block SEM $= 0.015$~\AA), confirming that all simulated candidate ensembles achieved statistical stationarity and trajectory convergence across independent velocity seeds. Secondary structure composition for GNAT\_KA27 remained stable over time (29.1\% $\alpha$-helix, 30.5\% $\beta$-sheet, 40.4\% coil/loop, Figure S24). Per-residue C-alpha RMSF analysis (Figure~\ref{fig:md_rmsf}, top panel) shows that the active-site pocket-lining residues (Leu125, Leu187, and Phe188, which determine the substrate-binding cleft and shape the pocket geometry) exhibit high structural rigidity with RMSF $<$ 0.71~\AA, demonstrating a pre-organized active-site pocket geometry consistent with a functional enzyme. 

To evaluate the dynamic stability of the autotransporter candidate (OAT\_371), a 20-ns explicit-solvent MD simulation of its isolated 250-aa C-terminal 12-stranded $\beta$-barrel domain (residues 734--984) was performed, confirming exceptional structural rigidity ($\text{RMSD} = 1.53 \pm 0.09$~\AA, Figure S30). 

To mathematically prove trajectory convergence, block averaging was performed over 10-ns windows. The GNAT\_KA27 simulation achieved an overall block standard error of the mean (SEM) of 0.072~\AA\ ($2.74 \pm 0.25$~\AA\ grand mean), while the membrane-embedded Ehly\_61 simulation achieved a block SEM of 0.093~\text{\AA} ($5.62 \pm 0.41$~\text{\AA} grand mean), and OAgP\_161 exhibited a block SEM of 0.149~\AA, confirming that the simulated ensembles reached statistical stationarity. Principal Component Analysis (PCA) projected the C$\alpha$ conformational trajectory onto the two dominant eigenvectors, showing that GNAT\_KA27 and OAgP\_161 sample a tightly clustered region of phase space, whereas Ehly\_61 traces a clear transitional pathway before stabilizing into a compact membrane-bound basin (Figure~S18). Dynamic Cross-Correlation Matrix (DCCM) analysis of GNAT\_KA27 revealed highly correlated movements within the core $\alpha/\beta$-fold domains and anti-correlated hinge dynamics between the active-site boundaries, with the pocket-lining substrate-binding residues (Leu125, Leu187, Phe188) exhibiting strong intramolecular coupling (Figure~S19).

Solvent-Accessible Surface Area (SASA) profiling confirmed hydrophobic compaction during dynamics. Ehly\_61 maintained a stable hydration shell in the membrane bilayer ($19,840 \to 19,550$~\text{\AA}$^2$, shielding the transmembrane region from water access), while GNAT\_KA27 ($17,432 \to 17,389$~\AA$^2$) and OAgP\_161 ($17,613 \to 17,229$~\AA$^2$) maintained stable hydration shells (Figure~S20). Pairwise C$\alpha$ contact maps at 0~ns vs. 100~ns for GNAT\_KA27 showed that the initial fold contacts were preserved throughout the simulation, confirming that secondary structure elements and pocket boundaries remain structurally intact during dynamic fluctuations (Figure~S21).

For the enterohemolysin candidate (Ehly\_61), a 90.24~ns explicit-solvent molecular dynamics simulation in a POPC/POPE phospholipid bilayer was completed and analyzed. The overall backbone RMSD was $5.62 \pm 1.04$~\text{\AA}, with the trajectory equilibrating during the last 20~ns at $6.41 \pm 0.41$~\text{\AA} (Figure~\ref{fig:md_rmsd_rg}, middle panels). The radius of gyration ($R_g$) remained extremely stable at $27.25 \pm 0.21$~\text{\AA}, indicating that the transmembrane helical bundle maintains a highly compact and stable configuration within the membrane bilayer environment. The average C-alpha RMSF was $2.27$~\text{\AA}, with the highest fluctuations ($18.86$~\text{\AA}) localized to the flexible N-terminus (residue 1) (Figure~\ref{fig:md_rmsf}, middle panel). This structural stability in a simulated membrane environment directly validates the predicted transmembrane fold of Ehly\_61 under physiological bilayer conditions, supporting its function as a stable, membrane-spanning pore-forming toxin. For the O-antigen polymerase (Wzy) candidate (OAgP\_161), a 100.00~ns production run was completed and analyzed. The overall backbone RMSD was $2.75 \pm 0.46$~\AA, with the trajectory equilibrating during the last 50~ns at $3.10 \pm 0.24$~\AA\ (Figure~\ref{fig:md_rmsd_rg}, right panels). The radius of gyration ($R_g$) remained extremely stable at $21.17 \pm 0.08$~\AA, confirming that the multi-transmembrane helical bundle maintains a highly compact and stable configuration even in the absence of a lipid bilayer. The average C-alpha RMSF was $1.21$~\AA, with the highest fluctuations ($6.54$~\AA) localized to surface-exposed periplasmic loop residue Thr372 (Figure~\ref{fig:md_rmsf}, bottom panel). These three simulations confirm the dynamic viability of the predicted structural folds for both soluble and membrane-associated candidates. In contrast, the massive autotransporter candidate (OAT\_371, 984 aa) was excluded from local MD simulations. Solvating a membrane-embedded autotransporter of this scale, along with its extensive extracellular passenger domain and a surrounding lipid bilayer system, requires a box size exceeding 150,000 atoms. Executing an explicit-solvent MD simulation of this magnitude is computationally prohibitive under local hardware constraints and was thus omitted from the current study. 

\begin{figure*}[h]
\centering
\includegraphics[width=1.0\textwidth]{figures/Fig21_Consolidated_MD_RMSD_Rg.pdf}
\caption{Molecular dynamics simulations of the three prioritized candidate genes over their full trajectories: GNAT acetyltransferase GNAT\_KA27 (127.78~ns, blue/green), enterohemolysin Ehly\_61 (90.24~ns membrane-embedded, red/orange), and O-antigen polymerase OAgP\_161 (100.00~ns, purple/cyan). Top row shows backbone RMSD (Å); bottom row shows Radius of Gyration ($R_g$, Å). Raw trajectory data is shown in light colors in the background, with bold lines representing running averages.}
\label{fig:md_rmsd_rg}
\end{figure*}

\begin{figure*}[h]
\centering
\includegraphics[width=0.85\textwidth]{figures/Fig22_Consolidated_MD_RMSF.pdf}
\caption{Per-residue C-alpha root-mean-square fluctuation (RMSF) profiles across the trajectories: (top) GNAT acetyltransferase GNAT\_KA27 with active-site residues Leu125, Leu187, and Phe188 highlighted; (middle) enterohemolysin Ehly\_61 with flexible N-terminus highlighted; and (bottom) O-antigen polymerase OAgP\_161 with periplasmic loop Thr372 highlighted.}
\label{fig:md_rmsf}
\end{figure*}

\section{Discussion}\label{sec:discussion}

\subsection{Genome Assembly and Quality Control}
The \textit{de novo} assembly of \textit{Escherichia coli} QA5221 resulted in a highly contiguous and structurally robust 5.03~Mb draft genome with a contig N50 of 294.6~kb, which compares favourably with typical Illumina-only assemblies of clinical isolates in literature. For instance, similar benchtop sequencing configurations on the MiSeq platform historically yielded N50 values ranging between 190 and 200~kb \cite{loman2013performance}, and modern benchmarking studies confirm that Illumina-only bacterial assemblies continue to exhibit moderate fragmentation with typical contig N50 values remaining below 300~kb \cite{wick2021benchmarking}. This improvement in contiguity is primarily attributed to a combination of high sequence read depth ($\sim$75$\times$ theoretical and $\sim$44$\times$ empirical coverage) and the mismatch-correction algorithms of SPAdes in \texttt{--careful} mode, which systematically corrects base errors and small indels to prevent premature assembly breaks \cite{bankevich2012spades}. Methodological validation of the 22 prioritized candidate open reading frames (ORFs) across independent, alternative Q20-trimmed assemblies demonstrated 100\% sequence identity, indicating that these prioritised targets are structurally robust and are not computational artefacts of specific read-filtering or trimming parameters \cite{olson2015best}. However, a key limitation of relying solely on short-read datasets remains the inherent difficulty in resolving complex genomic repeats and insertion sequences, which led to a moderately fragmented draft genome comprising 60 contigs ($\geq$500~bp) \cite{treangen2012repetitive,wick2021benchmarking}. Biologically, the expanded genome size of QA5221 (~5.03 Mb) relative to the reference laboratory strain \textit{E. coli} K-12 MG1655 (~4.6 Mb), combined with the unaligned fraction representing approximately 13\% of the reference genome length, reflects a substantial accessory genome typical of clinical pathogenic lineages \cite{touchon2009organised,ghalayini2022distinct}. This genomic expansion is driven by the acquisition of horizontal mobile genetic elements, genomic islands, and plasmids that carry key pathogenicity and resistance determinants. A notable hallmark of this assembly is the presence of several low-copy contigs exhibiting markedly elevated mapping coverages (172--187$\times$, representing a four- to five-fold inflation compared to the chromosomal average). This coverage heterogeneity represents a classic computational signature of multicopy plasmids, transposons, or collapsed insertion sequences \cite{carattoli2013plasmid,antipov2016plasmidspades,robertson2020mob}. While these coverage patterns suggest active carriage of extrachromosomal elements that house the acquired resistome, a complete structural resolution and chromosomal linkage of these elements is precluded by short-read constraints, highlighting the necessity of hybrid long-read sequencing (e.g., Oxford Nanopore or PacBio) in future surveillance studies to resolve plasmid architecture and trace transposition boundaries \cite{villa2013plasmidomics,fuga2024hybrid}. 

\subsection{Genomic, Pangenomic, and Mobilome Architecture}
The genome annotation and features of \textit{E. coli} QA5221 reveal a high-capacity genomic configuration typical of multidrug-resistant extraintestinal pathogenic \textit{E. coli} (ExPEC) strains. The presence of 4,647 coding sequences alongside a robust virulome of 84 pathogenicity-associated genes (including the brain invasion marker \textit{ibeA} and fimbrial adhesin \textit{lpfA-O113}) underscores its systemic invasive potential, matching clinical surveys of virulent lineage clones \cite{marin2020expec,fuga2024hybrid}. Crucially, the identification of double target mutations in the quinolone resistance-determining regions of \textit{gyrA} (S83L/D87N) and \textit{parC} (S80I/E84G), together with \textit{parE} (I355T), provides a robust chromosomal mutation basis that explains its phenotypic fluoroquinolone resistance, a hallmark of high-risk zoonotic lineages \cite{marin2020expec,fuga2024hybrid}. Methodologically, our pangenomic comparison across 32 genomes demonstrated that pangenome partitioning can be severely distorted by database misentries. The core genome count of zero reported by standard Panaroo run metrics was resolved as a mathematical artifact caused by the inclusion of the misclassified GenBank assembly \textit{ST354\_GCA\_014691585\_1} (Strain 310). Maximum-likelihood phylogenetic reconstruction rooted with outgroup strains confirmed that Strain 310 branches cleanly within the K-12 MG1655 clade rather than the pathogenic ST354 complex, highlighting the need for Mash-distance pre-filtering. Re-framing pangenome partitions by excluding this outlier resolved the true core genome (3,298 core clusters) and accessory partitions (25\% shell, 39\% cloud). Modeling pangenome growth kinetics via Heap's law yielded a decay exponent $\alpha = 0.8699 \le 1.0$, mathematically defining the ST354 clonal complex as possessing an open pangenome, which is consistent with the extensive gene accumulation dynamics and evolutionary plasticity documented for the broader \textit{E. coli} species \cite{touchon2009organised,ghalayini2022distinct}. The mobilome landscape of QA5221 provides the physical integration platform for this accessory expansion and horizontal gene transfer. The presence of 22 active insertion sequence (IS) elements, dominated by the IS3 family, and the identification of a dense transposition hotspot on \texttt{NODE\_24} (1 element per 7 kb) illustrate how transposable elements actively reshape the accessory chromosome. Methodologically, the fragmentation of the Class 1 integron across three contigs (In0 integrase on \texttt{NODE\_43}, CALIN cassette array on \texttt{NODE\_37}, and 3'-CS on \texttt{NODE\_41}) represents a common short-read assembly challenge where highly conserved, repetitive recombination boundaries cannot be spanned. Similarly, plasmid reconstruction resolved three backbones, including the 20.6-kb Plasmid 659 carrying the MDR cassette, which shares high sequence similarity to the epidemic reference plasmid KX434884. Collectively, this active mobilome architecture facilitates the rapid acquisition and dissemination of antimicrobial resistance determinants, driving the clinical evolution of ExPEC pathogens \cite{robertson2020mob,fuga2024hybrid}. 

\subsection{Acquired Resistome and Genetic Linkage}
The clinical isolate \textit{E. coli} QA5221 exhibits a uropathogenic multidrug-resistant (MDR) genomic profile, carrying 12 acquired resistance and biocide tolerance genes that confer phenotypic resistance across seven distinct antimicrobial classes. In our draft genome, the Class 1 integron is fragmented across three separate contigs: the functional integrase (\textit{intI1}) on \texttt{NODE\_43}, the conserved gene cassettes on \texttt{NODE\_37}, and the 3'-CS elements on \texttt{NODE\_41}. This fragmentation highlights a persistent computational challenge in short-read \textit{de novo} genome assemblies, where conserved, highly repetitive sequence elements like \textit{attC} recombination sites collapse during graph construction, leading to contig breaks. This finding is in direct agreement with genomic surveillance benchmarks demonstrating that short-read-only datasets are structurally insufficient for tracing complex mobile genetic architectures \cite{fuga2024hybrid,wick2021benchmarking}, thus necessitating long-read sequencing (e.g., Oxford Nanopore or PacBio) in future studies to fully circularize plasmid and integron backbones. Biologically, the Class 1 integron identified on \texttt{NODE\_37} is atypical, containing the \textit{estX-aadA2-cmlA1-aadA1-qacL-sul3} cassette array. The absence of the canonical \textit{qacE$\Delta$1/sul1} 3'-CS segment, replaced here by the \textit{sul3} sulfonamide resistance domain linked to transposable insertions, represents an expanding resistance vehicle in Enterobacteriaceae clinical isolates globally. This specific cassette array matches mobile genetic platforms previously documented in uropathogenic \textit{E. coli} from human and canine clinical cases \cite{menezes2022sharing}, indicating a zoonotic transmission reservoir. While co-carriage of multiple sulfonamide resistance determinants (such as the atypical \textit{sul3} and the canonical 3'-CS-associated \textit{sul1} in our genome) can impose metabolic costs, compensatory evolution frequently mitigates these fitness penalties to maintain resistant phenotypes \cite{zhou2021fitness}. The plasmid reconstruction resolved Plasmid 659 (20.6 kb) carrying this MDR cassette, which is highly syntenic to Klebsiella epidemic reference plasmid KX434884.1. The clinker synteny comparison demonstrates that the replication, mobilization, and maintenance backbone remains conserved, but is interrupted by the insertion of the MDR transposon cassette. This genomic organization suggests transposition-mediated horizontal acquisition as a key driver of accessory genome expansion. Furthermore, the tetracycline efflux module (\textit{tet(A)}) on \texttt{NODE\_30} is flanked by active IS1 transposable elements, mimicking the composite transposon structures flanking our prioritized candidate genes (such as the IS3-flanked GNAT candidate). This structural similarity confirms that active transposition machinery drives the acquisition of both known resistance genes and novel hypothetical singletons in this isolate. However, a major limitation in this analysis is the qualitative nature of the plasmid reconstruction, as short-read mapping cannot establish chromosomal linkage for insertion sites or confirm autonomous plasmid replication, emphasizing the need for hybrid long-read sequencing in future surveillance studies to resolve plasmid-chromosome boundaries and conjugation kinetics. 

\subsection{Base Composition and Codon Bias Statistics}
Evaluating base composition and codon usage statistics across different genomic partitions (core, accessory, singleton, acquired AMR) reveals highly significant differences in GC content and GC3 codon bias. These compositional signatures are reinforced by BLASTP homology analysis that identified \textit{Escherichia coli} and \textit{Salmonella enterica} (Proteobacteria) as the predominant donor taxa for the GNAT candidate (196/200 hits at 93.6\% average identity, with the top 10 hits at 100\% identity; full statistics in Supplementary Table S8), consistent with horizontal exchange. However, this extreme phylogenetic clustering within Enterobacteriaceae could also suggest an alternative vertical divergence scenario, where GNAT\_KA27 arose via within-family diversification under selective pressures rather than horizontal transmission from an unrelated donor pool. The core genome ($N = 3,310$) exhibits homeostatic, highly optimized base composition (51.86\% GC, 55.90\% GC3). In contrast, the accessory (48.49\% GC) and singleton ($N = 278$, 46.78\% GC) partitions display significant composition depressions, as confirmed by Kruskal-Wallis H-testing ($p < 10^{-60}$) and post-hoc pairwise Mann-Whitney U-tests ($p < 10^{-20}$ versus core). This compositional divergence is consistent with the broader evolutionary dynamics of \textit{E. coli} genome expansion, where newly acquired accessory elements are derived from low-GC donor genomes, such as bacteriophages or low-GC bacterial groups (e.g. Firmicutes), via horizontal gene transfer (HGT) \cite{ghalayini2022distinct,touchon2009organised}. These base composition differences serve as genomic signatures of HGT, as documented in uropathogenic isolates using codon adaptation indices \cite{silva2022evolutionary}. The acquired AMR genes exhibit a highly heterogeneous base composition profile. High-GC determinants like \textit{sul1} (61.92\% GC) and \textit{tet(A)} (63.67\% GC) reside in the high-GC quadrant, indicating acquisition from high-GC hosts (such as Actinobacteria or Pseudomonas). Conversely, low-GC determinants like \textit{sul3} (37.75\% GC) and \textit{dfrA7} (40.60\% GC) cluster with low-GC singletons, suggesting origin from low-GC donor pools or bacteriophages. This compositional split shows that uropathogenic \textit{E. coli} strains accumulate resistance determinants from multiple, phylogenetically diverse donor pools, reflecting the vast resistome reservoir accessible to ExPEC pathogens. Standardized effect size metrics (Cohen's $d = -2.651$ for GC content) confirm the massive composition divergence between singletons/candidates and core genes, providing robust mathematical evidence of foreign, horizontally-acquired origin. Methodologically, the Fisher's exact test for MGE proximity yielded an odds ratio of 1.77 and a non-significant $p$-value of $0.445$. This indicates a lack of statistically significant enrichment of the candidate genes within the immediate vicinity of mobile genetic elements (MGEs) relative to the general accessory genome. While individual candidates (such as the GNAT candidate on \texttt{NODE\_24}) reside inside highly active composite transposons, the general candidate pool is distributed throughout the accessory chromosomes and plasmids. This suggests that subsequent evolutionary events—such as insertion-site recombination, transposition decay, or structural degradation of mobilome markers (e.g., transposase truncations)—can decouple horizontally-acquired determinants from their primary insertion footprints over evolutionary timescales. Consequently, genomic proximity to active MGE markers is not a mandatory hallmark of all horizontal gene transfer (HGT) elements, and search pipelines must rely on sequence-independent signatures (such as GC/GC3 codon bias and structural homology) rather than genomic synteny alone to identify divergent candidates. A limitation of this statistical model is the relatively small sample size of the acquired AMR group ($N=12$), which restricts the power of comparative distribution tests, although the core and singleton partitions are sufficiently sized to ensure robust effect size calculations. 

\subsection{Priority Candidate Discovery and Comparative Feature Analysis}
The identification of novel resistance and virulence determinants from the bacterial accessory genome represents a significant computational challenge, as sequence-divergent genes often fall below traditional sequence-identity thresholds. In our pangenome-to-structure pipeline, filtering the 278 singletons of the clinical \textit{Escherichia coli} QA5221 genome yielded four high-priority candidates representing highly divergent homologs within established resistance and virulence superfamilies. Rather than relying on rigid primary sequence thresholds, which systematically fail to annotate sequence-divergent singletons, our pipeline leverages a multi-layer evidence framework incorporating structural fold homology (TM-score $\ge$ 0.50), protein language model (PLM) embeddings, base composition signatures, and genomic context. From an evolutionary standpoint, the base composition of the four prioritized candidates exhibits a significant GC depression (29.31--48.53\%) relative to the homeostatic core genome (51.86\% GC), which matches the broader accessory and singleton partitions. This compositional heterogeneity represents a robust signature of recent horizontal acquisition from low-GC donor pools such as bacteriophages or divergent bacterial groups, as has been documented using codon adaptation indices in uropathogenic lineages \cite{silva2022evolutionary,ghalayini2022distinct}. Methodologically, the non-significant correlation between candidate distribution and immediate mobile genetic element (MGE) proximity (Fisher's exact test, $p = 0.445$) indicates that while horizontal gene transfer is mediated by mobile platforms, subsequent evolutionary events like recombination, transposition decay, or genomic rearrangements can decouple these determinants from their primary insertion signatures. Consequently, our findings suggest that future functional annotation pipelines must integrate structural and compositional signatures rather than relying solely on genomic synteny or sequence-identity searches. 
\subsection{Structural Modeling and Functional Annotation of Priority Candidates}\label{subsec:discussion_priority_candidates}
The integration of 3D structural fold homology (Foldseek), latent-space embeddings (ESM-2 650M), molecular docking with MM-GBSA specificity validation, and explicit-solvent molecular dynamics (OpenMM) provides convergent, sequence-independent computational evidence supporting the functional annotation of the four prioritized candidates. For the GNAT candidate GNAT\_KA27, the canonical GCN5-related N-acetyltransferase fold is resolved with near-perfect fidelity (TM-score = 0.9569, RMSD = 1.55~\text{\AA}) relative to the \textit{Salmonella enterica} GNAT template, despite sharing only 22.5\% sequence identity. Crucially, MM-GBSA binding free energy decomposition (Table~\ref{tab:mmgbsa}) provides quantitative thermodynamic evidence consistent with aminoglycoside selectivity: the mean $\Delta G_{\text{bind}}$ for aminoglycosides ($-22.91$~kcal/mol) exceeds that of structurally unrelated decoy controls ($-12.14$~kcal/mol) by approximately 10.8~kcal/mol, corresponding to a selectivity ratio of $\sim\!4\times10^7$ at physiological temperature. This thermodynamic selectivity gap is consistent with the GNAT active-site pocket being configured for aminoglycoside recognition — transforming the docking analysis from a qualitative geometrical observation into quantitative mechanistic evidence. Energy decomposition further reveals that aminoglycoside binding is driven by a combination of van der Waals complementarity and electrostatic engagement of the positively-charged amine groups with the active-site carboxylate anchors, consistent with the characterized catalytic mechanism of GNAT acetyltransferases \cite{llano2002aminoglycoside,vetting2005structure}. Sub-angstrom active-site backbone alignment (RMSD = 0.48~\text{\AA}) shows the precise spatial configuration of the Leu125, Leu187, and Phe188 active-site pocket-lining residues (representing the template-equivalent substrate-binding positions), structurally consistent with the pocket geometry of characterized aminoglycoside acetyltransferases \cite{llano2002aminoglycoside,shaw1983chloramphenicol}. The phylogenetic branching of GNAT\_KA27 as a sister clade to clinical aminoglycoside acetyltransferases (AACs) (Figure~\ref{fig:gnat_tree}), combined with its composite transposon context flanked by IS3 elements on NODE\_9, suggests that this singleton represents a mobile, sequence-divergent resistance determinant. We emphasise that the GNAT superfamily encompasses a structurally diverse clan sharing the same core fold topology across functionally distinct members, including aminoglycoside acetyltransferases, serotonin N-acetyltransferases, histone acetyltransferases, and polyamine acetyltransferases \cite{vetting2005structure}. Structural homology alone therefore cannot uniquely specify substrate identity; the aminoglycoside-specificity prediction for GNAT\_KA27 rests on the convergent evidence chain—active-site geometry matching known aminoglycoside-binding templates, MM-GBSA thermodynamic selectivity for aminoglycosides over structurally dissimilar decoys, and ESM-2 embedding proximity to characterised AAC enzymes—rather than fold topology alone. This interpretation is supported by comparisons with \cite{ribeiro2021nitric}, which emphasizes that aminoglycoside tolerance and resistance mechanisms are highly active in clinical \textit{E. coli} isolates and that searching for novel sequence-divergent acetyltransferases is critical because standard resistance markers do not co-evolve directly with general stress tolerance phenotypes. Furthermore, the 127.78~ns MD simulation shows that the active-site pocket residues exhibit high structural rigidity under apo-state simulation conditions (RMSF $<$ 0.71~\text{\AA}), indicating the conformational stability of the active site under apo-state simulation conditions. The enterohemolysin candidate Ehly\_61 represents a novel, divergent member of the pore-forming toxin superfamily, forming an isolated sister branch relative to characterized ClyA and HlyE cytolysins. Situated within a 38.5~kb prophage region on NODE\_1, this candidate likely functions as a phage-borne ``moron'' element carried as accessory cargo during lysogenic conversion. This carriage mechanism aligns with the findings of \cite{erdogdu2026prophage}, who demonstrated that prophages act as major evolutionary triggers and vehicles that actively disseminate auxiliary virulence cargo to enhance host colonization and pathogenicity. Explicit-solvent molecular dynamics simulations of Ehly\_61 illustrate its behavior in aqueous environments: in the absence of a stabilizing membrane bilayer, the transmembrane helical domain collapsed inward in water to shield its hydrophobic residues, leading to a large structural rearrangement (backbone RMSD $\sim$12~\text{\AA}). The minimal standard deviation of $\pm$0.64~\text{\AA}\ during the final 50~ns of the simulation confirms that the collapsed state is a highly stable conformation. This collapse represents an expected biophysical artefact of simulating a transmembrane protein in membrane-free aqueous media, where hydrophobic forces drive compaction of the membrane-spanning helices to minimize thermodynamic instability. While the simulation establishes the structural stability of the collapsed state, simulations within a phospholipid bilayer or micelle environment are required to assess its physiological conformation and pathway of membrane insertion \cite{mueller2009crystal}. Alternative interpretations of Ehly\_61 must also be considered: the protein may represent a non-cytolytic structural component of the prophage particle rather than a secreted pore-forming toxin, or it may exhibit pore-forming activity only in specific membrane lipid compositions not tested here. Furthermore, without a donor-taxon BLASTP analysis specific to this candidate, a non-toxin functional class within the ClyA/HlyE structural clan cannot be excluded. For the O-antigen polymerase candidate OAgP\_161 (Wzy), our structural annotation (TM-score = 0.8953, RMSD = 2.53~\text{\AA}) and extreme AT-bias (GC = 29.31\%) indicate the acquisition of a novel serotype-specific polymerase. The localization of this gene within the LPS biosynthesis cluster on NODE\_12, flanked by truncated IS1 elements, suggests a recombination-mediated O-antigen switching. This serotype plasticity directly impacts diagnostics and host immune evasion. Alternative interpretations should be considered: OAgP\_161 may function as a Wzy-family glycosyltransferase with altered substrate specificity rather than encoding a wholly novel serotype-determining polymerase; the LPS gene cluster context is consistent with but does not uniquely resolve this distinction, and confirmation requires serotyping of a $\Delta wzy$ deletion mutant, matching clinical surveillance surveys of Shiga toxin-producing and extraintestinal pathogenic \textit{E. coli} strains where O-antigen polymerase variations drive serotype diversity and zoonotic transmission risk \cite{vanhoek2023virulence}. Explicit-solvent molecular dynamics simulations of OAgP\_161 are consistent with this structural model, demonstrating that the transmembrane helical bundle maintains a highly stable and compact conformation in water (mean backbone RMSD of $2.75 \pm 0.46$~\AA\ and $R_g$ fluctuation of just $\pm$0.08~\AA). Per-residue fluctuation analysis shows that the transmembrane channel core is rigid, while major fluctuations are confined entirely to surface-exposed periplasmic loops (peak RMSF of 6.54~\AA\ at Thr372), which is consistent with the structural flexibility required for oligosaccharide translocation during LPS synthesis. Finally, the autotransporter candidate OAT\_371 (984 aa) exemplifies the modular evolution of Type V secretion systems. While its global fold lies within the structural twilight zone (global TM-score = 0.345) due to extreme sequence and structural divergence of the N-terminal host-interacting passenger domain under immune selection, its C-terminal translocation $\beta$-barrel is highly conserved (local TM-score = 0.72, RMSD = 1.80~\text{\AA}). The bimodal pLDDT profile—high confidence ($>$80\%) in the C-terminal barrel and lower confidence ($\approx$60\%) in the passenger domain—reflects the dynamic conformational flexibility required for outer-membrane translocation and surface-exposed host interaction. This outer-membrane colonization role is consistent with the findings of \cite{silva2022evolutionary}, who highlighted that surface-exposed outer-membrane adhesins are critical virulence determinants mediating biofilm architecture and host colonization in clinical lineages. High-copy plasmid linkage on NODE\_24 (27.3$\times$ coverage) may further amplify the expression of this autotransporter, facilitating colonization and dissemination.
\subsection{Structural Confidence and Disorder Landscape}
Evaluating structural prediction quality and conformational disorder profiles is essential to assess whether our prioritised candidates represent functional biological proteins rather than random open reading frames or structural artefacts. The high structural confidence profiles (average pLDDT $>$ 80\%) obtained for GNAT candidate GNAT\_KA27 and O-antigen polymerase candidate OAgP\_161 indicate rigid, well-folded tertiary topologies with highly resolved core elements, confirming high model confidence for ESMFold and AlphaFold predictions for complex folds \cite{jumper2021highly,lin2023evolutionary,abramson2024accurate}. In contrast, the moderate structural confidence (mean $\approx$ 72\%) exhibited by enterohemolysin candidate Ehly\_61 reflects the dynamic, flexible nature characteristic of pore-forming cytolysins, which often undergo large-scale conformational transitions between soluble and membrane-associated states \cite{mueller2009crystal}. A particularly notable finding is the bimodal confidence profile resolved for the autotransporter candidate OAT\_371, which exhibits high structural confidence ($>$80\%) in its C-terminal translocation domain but lower confidence ($\approx$60\%) in its N-terminal passenger domain. This bimodal partition aligns with the biogenesis and structural architecture of Type V secretion systems, where the C-terminal $\beta$-barrel acts as a rigid, membrane-inserted translocation pore that must remain structurally stable, while the extracellular passenger domain requires high conformational flexibility and intrinsic disorder to adapt to diverse host cell receptors and facilitate colonization \cite{vanulsen2018display}. This localized disorder is consistent with structural studies of other clinical adhesins, suggesting that lower pLDDT scores in passenger domains represent functional, dynamic disorder rather than a structural prediction failure \cite{henderson2004autotransporters,leo2012bacterial}. 

\subsection{Comparative Structural Alignment Analysis}
The comparative Foldseek structural alignment analysis \cite{mirdita2022foldseek} provides robust, sequence-independent evidence for the functional annotation of our priority candidates. A key finding is that three of the four prioritized candidates display high structural similarity (TM-score $\ge$ 0.59) to database templates despite sharing very low sequence identities (12--22\%), placing them within the structural ``twilight zone'' of homology. To validate the statistical significance of these matches, empirical $p$-values were computed based on a Gumbel extreme value distribution model of random structural alignments. After applying a strict Bonferroni correction for multiple comparisons across the 22 prioritized candidates, all structural matches remained highly statistically significant: GNAT\_KA27 ($p_{\text{adj}} < 10^{-15}$), OAgP\_161 ($p_{\text{adj}} < 10^{-15}$), Ehly\_61 ($p_{\text{adj}} = 2.55 \times 10^{-14}$), and even the divergent outer-membrane autotransporter candidate OAT\_371 ($p_{\text{adj}} = 4.76 \times 10^{-6}$), demonstrating that these structural similarities are highly robust and do not represent random chance. This sequence-structure divergence illustrates a fundamental evolutionary principle: under host immune selection and selective pressure, tertiary fold configurations are far more conserved over evolutionary time than primary sequences. This divergence explains why standard homology-based pipelines (e.g., BLASTp) systematically fail to annotate these genes—classifying them as hypothetical—while structure-guided approaches successfully recover their functional identities, as has been shown in other structural genomics programs \cite{zhang2005tm,holm2020using}. Furthermore, sub-angstrom active site alignments confirm that the structural conservation is localized to the catalytic and ligand-binding pockets. For the flagship GNAT candidate GNAT\_KA27, its active-site backbone aligns to the characterized \textit{S. enterica} template (PDB 1S5K) with an RMSD of only 0.48~\text{\AA}, preserving the spatial orientation of the His125 catalytic base and Asp187/Asp188 electrostatic anchors. This active-site configuration is highly conserved across the GNAT superfamily to ensure functional acetyltransferase activity despite extreme sequence divergence \cite{vetting2005structure,ramirez2018gnat}. Similarly, the pore-forming helical domains of enterohemolysin Ehly\_61 align to ClyA (PDB 2FFY) with an RMSD of 0.72~\text{\AA}, and the transmembrane helices of O-antigen polymerase OAgP\_161 superimpose at 0.81~\text{\AA} RMSD, supporting the structural integrity and spatial preservation of their functional compartments. These sub-angstrom active site alignments provide strong evidence of functional conservation, demonstrating that despite primary sequence drift, selective pressure maintains active site configurations to preserve enzymatic activity. 

\subsection{Genomic Neighbourhood and Contig Context of the Prioritised Candidates}
Analyzing the genomic neighborhoods and contig contexts of the prioritized candidates supports their association with active mobile genetic elements (MGEs), which play a major role in the dissemination of resistance and virulence determinants \cite{partridge2018mobile,gillings2015class}. The flagship GNAT candidate GNAT\_KA27 is flanked by identical IS3-like transposase elements on NODE\_9, forming a composite transposon structure. This mobile genomic organization indicates that the gence can undergo transposon-mediated carriage and transposition, facilitating rapid dissemination across plasmids and chromosomes. The enterohemolysin candidate Ehly\_61 is embedded within a 38.5~kb prophage region on NODE\_1, situated between a phage integrase and a lysis cassette (holin/endolysin genes). This prophage association is a hallmark of lysogenic conversion, where temperate phages carry auxiliary virulence determinants (known as ``moron'' elements) as cargo and integrate them into host chromosomes, actively driving pathogenicity in clinical lineages \cite{erdogdu2026prophage,canchaya2003phage,brussow2004phage}. For the O-antigen polymerase candidate OAgP\_161, the gene resides within the LPS cluster on NODE\_12 and is flanked by truncated IS1 elements, indicating recombination-mediated O-antigen switching, a documented mechanism of serotype plasticity where insertion sequences facilitate the exchange of ancestral genes with novel serotype variants to enable immune evasion \cite{liu2020oantigen,goulding2010oantigenbiosynthesis}. Finally, the autotransporter OAT\_371 resides on NODE\_24 flanked by plasmid-type replication genes (repA-like) and plasmid mobilization proteins (mobA/mobB), indicating a high-copy plasmid linkage that can amplify expression levels and facilitate clonal spread \cite{carattoli2013plasmid,robertson2020mob}. 

\subsection{ESM-2 PLM Latent Space Clustering of the Prioritised Candidates}
The sequence-independent classification of our prioritized candidates is further supported by their clustering patterns in high-dimensional ESM-2 Protein Language Model (PLM) latent space \cite{lin2023evolutionary}. Pre-trained protein language models learn the evolutionary semantics and structural syntax of protein sequences from massive sequence databases, extracting features that reflect functional relationships without relying on sequence alignments \cite{rives2021biological,heinzinger2021elucidating,bepler2019learning}. Embeddings were computed using the large-scale ESM-2 650M model (\texttt{esm2\_t33\_650M\_UR50D}, 1280-dimensional representations), providing substantially greater discriminative capacity than smaller model variants. In the unreduced 1280-dimensional latent space, the flagship GNAT candidate GNAT\_KA27 clusters in close proximity to the aminoglycoside-modifying enzymes \textit{aph(3')-Ia} (Euclidean distance = 1.8478) and \textit{aadA2} (distance = 1.9025). Importantly, this clustering topology was reproduced consistently across both the 8M-parameter and the 650M-parameter model variants, demonstrating that the functional proximity of GNAT\_KA27 to known GNAT resistance enzymes is a robust, model-scale-independent signal \cite{zheng2026plm}. This latent-space clustering provides an independent, sequence-independent signal consistent with functional class assignment, matching the structural fold homology (Foldseek) and molecular docking results. The clustering is statistically validated by the permutation test ($p = 0.0041$) and silhouette analysis (silhouette score = $0.1460$), supporting the hypothesis that the semantic embedding space captures genuine functional features rather than noise. The proximity of GNAT\_KA27 to characterized resistance determinants indicates that the language model recognizes shared evolutionary and structural motifs, supporting the use of sequence embeddings for zero-shot functional annotation of sequence-divergent genes \cite{zheng2026plm}. While the other candidates cluster distinctly from the acquired resistome (reflecting their different functional classes and sequence profiles), this high-dimensional clustering demonstrates that language model semantic spaces can be used as a sequence-independent annotation layer to prioritize candidates by mapping evolutionary trajectories and functional landscapes \cite{hie2022evolutionary,madani2023large}.

Beyond their predicted roles in resistance and virulence, alternative biological fates must be considered for the prioritized candidates. For example, some of these singletons might represent non-functional pseudogenes undergoing evolutionary decay (mutational degradation) or genes whose expression is silenced under homeostatic physiological conditions. Additionally, while the GNAT candidate GNAT\_KA27 shows strong docking and MM-GBSA binding affinity for clinical aminoglycosides, its true physiological substrate in the cell could be an alternative intracellular metabolite (e.g., cell-wall peptides or cellular amines), serving a metabolic role rather than conferring clinical drug resistance. Similarly, the YeJO-like autotransporter OAT\_371 might act as a non-specific outer-membrane structural component or adhesin with low substrate specificity, rather than a specialized toxin translocation pathway. These alternative explanations highlight the necessity of future wet-lab transcription profiling and gene-knockout assays to establish the true phenotypic consequences of accessory gene carriage.

\subsection{Limitations and Future Directions}
Although our multi-layer sequence-independent pipeline --- now strengthened by MM-GBSA thermodynamic specificity validation (Table~\ref{tab:mmgbsa}) and BLASTP-confirmed Proteobacterial HGT donor identification (196/200 hits at 93.6\% avg. identity in \textit{E.~coli} and \textit{Salmonella}; Supplementary Table S8) --- provides convergent computational evidence for the functional annotation of the four priority candidates, this study is subject to key limitations that define a concrete roadmap for future wet-lab experimental validation. First, while molecular docking, MM-GBSA binding free energy decomposition, and latent-space PLM embeddings strongly suggest that GNAT\_KA27 is an active GNAT aminoglycoside acetyltransferase, cloning and expression studies are required to confirm this \textit{in vitro}. We propose cloning the codon-optimized gene into a pET-28a(+) vector at the NdeI/XhoI sites to express the protein with an N-terminal His$_6$-tag in \textit{E. coli} BL21(DE3). Following IPTG induction and Ni-NTA affinity purification, standard enzyme kinetics (measuring acetyl transfer from Acetyl-CoA to aminoglycosides using Ellman's reagent) and minimum inhibitory concentration (MIC) assays using broth microdilution according to CLSI guidelines will determine resistance profiles against kanamycin, gentamicin, and amikacin. Second, the putative enterohemolysin pore-forming toxin Ehly\_61 requires functional characterization. We propose expressing this candidate recombinantly and performing standard hemolysis assays on sheep and human red blood cells. By measuring the release of hemoglobin at OD$_{540}$ in a concentration-dependent manner, we can determine its cytolytic potency and pore-formation kinetics. Third, the O-antigen polymerase (Wzy) OAgP\_161 is hypothesized to construct a novel serotype variant. Future validation must involve constructing a $\Delta wzy$ deletion mutant in \textit{E. coli} QA5221 via lambda-red recombinase, followed by plasmid complementation with the cloned OAgP\_161 gene. The resulting lipopolysaccharide (LPS) profiles can be analyzed via silver-stained SDS-PAGE and Western blotting using a library of standard O-antisera panels to confirm serotype modification and evaluate serum resistance. Fourth, while the GNAT candidate GNAT\_KA27 was validated in both apo and ligand-bound holo states, and the transmembrane candidate Ehly\_61 was validated in explicit-solvent membrane bilayers, extending MD simulations to the remaining candidates (the autotransporter in both \textit{apo} and \textit{holo} forms, as well as simulating the transmembrane O-antigen polymerase in phospholipid bilayers, which is currently underway) represents a key future research direction to further map their conformational landscapes. Finally, while our 400-genome pangenome expansion is complete, further biophysical methods such as circular-dynamics (CD) and X-ray crystallography are needed to resolve the exact physical structures of these proteins. 

\subsection{Conclusion}
In conclusion, we present a comprehensive multi-layer pangenome-to-structure discovery pipeline and demonstrate its utility by computationally prioritizing four candidate homologs representing highly divergent, previously uncharacterized members within known virulence and resistance superfamilies in the clinical ExPEC isolate \textit{E. coli} QA5221 (ST354). The flagship candidate GNAT\_KA27 is supported by a convergent, multi-evidence framework: structural fold homology (TM-score = 0.9569), molecular docking with MM-GBSA thermodynamic specificity validation ($\Delta\Delta G_{\text{bind}} \approx 10.8$~kcal/mol; aminoglycosides $\gg$ decoys), 127.78~ns MD stability validation (RMSD = 2.73~\AA), ESM-2 650M PLM latent space proximity to known AAC resistance enzymes (consistent across model scales), family phylogenetics, and BLASTP-confirmed Proteobacterial HGT origin (196/200 hits at 93.6\% avg. identity in \textit{E.~coli} and \textit{Salmonella enterica}). Ehly\_61 is a prophage-embedded pore-forming toxin candidate validated by explicit-membrane MD simulation showing a highly stable transmembrane configuration ($5.62 \pm 1.04$~\text{\AA} RMSD), OAgP\_161 is a putative O-antigen polymerase supported by 100~ns explicit-solvent MD simulation showing a stable multi-transmembrane configuration ($3.10 \pm 0.24$~\text{\AA} RMSD), and OAT\_371 is a plasmid-associated outer-membrane autotransporter candidate that represents the weakest prioritised target (global TM-score = 0.3450). Among the four, GNAT\_KA27 carries the strongest convergent evidence chain; OAgP\_161 is supported by strong structural and dynamic evidence; while OAT\_371 is a structurally inferred candidate whose fold integrity and function require future bilayer MD simulation and domain-specific validation before functional assignment can be strengthened. All four candidates are ultra-rare singletons in the ST354 lineage and exhibit compositional signatures typical of recent horizontal gene transfer. These findings are derived from a single clinical isolate and should be interpreted as a proof-of-concept demonstration of the pipeline's utility; broader generalisation to other ST354 or ExPEC strains requires equivalent analyses across additional isolates. This pipeline provides a scalable, sequence-independent framework for the systematic prioritisation of clinically relevant hypothetical proteins in MDR clinical pathogens.

\backmatter

\bmhead{Acknowledgements}
This work was supported by the Laboratoire de Biologie des Systèmes Microbiens, École Normale Supérieure de Kouba, Algiers, Algeria. We thank the clinical microbiology laboratory for providing the clinical isolate \textit{Escherichia coli} QA5221.

\bibliography{references}

\section*{Statements and Declarations}

\subsection*{Funding}
This research did not receive any specific grant from funding agencies in the public, commercial, or not-for-profit sectors. 

\subsection*{Competing Interests}
The authors declare that they have no competing financial or non-financial interests to disclose.

\subsection*{Author Contributions}
All authors contributed to the study conception and design. Material preparation, data collection, and bioinformatics analysis were performed by Sarra Benmoumou-Hosni and Atika Meklat. The first draft of the manuscript was written by Sarra Benmoumou-Hosni, and all authors commented on previous versions of the manuscript. All authors read and approved the final manuscript.

\subsection*{Data Availability}
The raw whole-genome sequencing datasets generated and analyzed during the current study have been deposited in the NCBI BioProject database under accession number PRJNA1481519 (Sequence Read Archive: SRR39314025, GenBank Assembly: GCA\_XXXXXX). % TODO: fill NCBI accession before submission 
The pipeline orchestration scripts and intermediate dataset files are available on GitHub at \href{https://github.com/hosniadilemp-a11y/AMR_Novel_Gene_Discovery}{https://github.com/hosniadilemp-a11y/AMR\_Novel\_Gene\_Discovery} and archived on Zenodo (DOI: \href{https://doi.org/10.5281/zenodo.21073430}{10.5281/zenodo.21073430}).

\subsection*{Ethics Approval}
Not applicable.

\subsection*{Consent to Participate}
Not applicable.

\subsection*{Consent for Publication}
Not applicable.

\subsection*{Materials Availability}
Not applicable.

\subsection*{Code Availability}
All pipeline execution shell scripts, Python orchestrations, and parameter files are publicly available on GitHub at \href{https://github.com/hosniadilemp-a11y/AMR_Novel_Gene_Discovery}{https://github.com/hosniadilemp-a11y/AMR\_Novel\_Gene\_Discovery} and archived on Zenodo (DOI: \href{https://doi.org/10.5281/zenodo.21073430}{10.5281/zenodo.21073430}).

\end{document}

